% !TeX program = xelatex
% !TeX encoding = UTF-8
% 이전에 제공한 한국어 번역 PDF를 바탕으로 재구성한 편집 가능한 LaTeX 파일.
% 원래 번역 PDF를 제작한 소스 파일 자체를 복구한 것은 아닙니다.
% 원문 기준 커밋: f9e39354e4625e070c100a9ef2d55c77552e2076
% 외부 .bib 파일 없이 이 파일만으로 XeLaTeX 컴파일이 가능합니다.
% 컴파일: xelatex latent_art_bench_Korean_translation.tex (두 번 실행)
% 그림은 기존 번역본과 같이 원본 링크로 표시합니다.
% figures/ 아래에 해당 PDF를 넣으면 실제 그림이 자동으로 삽입됩니다.
\documentclass[11pt,a4paper]{article}
\usepackage[margin=25mm,headheight=15pt]{geometry}
\usepackage{iftex}
\ifXeTeX\else
  \PackageError{KoreanTranslation}{Please compile with XeLaTeX}{Select XeLaTeX as the compiler.}
\fi
\usepackage{fontspec}
\usepackage{kotex}
\IfFontExistsTF{TeX Gyre Termes}{\setmainfont{TeX Gyre Termes}}{\setmainfont{Liberation Serif}}
\IfFontExistsTF{TeX Gyre Heros}{\setsansfont{TeX Gyre Heros}}{\setsansfont{Liberation Sans}}
\IfFontExistsTF{TeX Gyre Cursor}{\setmonofont{TeX Gyre Cursor}[Scale=0.90]}{\setmonofont{Liberation Mono}[Scale=0.90]}
\IfFontExistsTF{Noto Serif CJK KR}{
  \setmainhangulfont{Noto Serif CJK KR}[AutoFakeSlant=0.15]
}{\IfFontExistsTF{NanumMyeongjo}{
  \setmainhangulfont{NanumMyeongjo}[BoldFont=NanumMyeongjo Bold,AutoFakeSlant=0.15]
}{\setmainhangulfont{UnBatang}[AutoFakeSlant=0.15]}}
\IfFontExistsTF{Noto Sans CJK KR}{
  \setsanshangulfont{Noto Sans CJK KR}
}{\IfFontExistsTF{NanumGothic}{\setsanshangulfont{NanumGothic}}{\setsanshangulfont{UnDotum}}}
\usepackage{amsmath,amssymb,graphicx,booktabs,tabularx,array}
\usepackage[round,authoryear]{natbib}
\usepackage[font=small,labelfont=bf]{caption}
\usepackage{float,placeins,fancyhdr,xcolor}
\usepackage{xurl}
\usepackage[unicode,colorlinks=true,allcolors=blue,breaklinks=true]{hyperref}
\hypersetup{
  linkcolor=black,citecolor=black,urlcolor=blue,
  pdftitle={텍스트-이미지 생성에서 화가 이름이 유발하는 반응: 공통 회화 효과를 넘어서 — 한국어 번역},
  pdfauthor={원문: 익명 저자 / 한국어 번역: ChatGPT},
  pdfsubject={한국어 번역 PDF를 바탕으로 재구성한 LaTeX 편집본; 기준 커밋 f9e39354e462}
}
\setlength{\parindent}{1em}
\setlength{\parskip}{3pt}
\setlength{\emergencystretch}{3em}
\linespread{1.12}
\clubpenalty=10000
\widowpenalty=10000
\renewcommand{\topfraction}{0.9}
\renewcommand{\bottomfraction}{0.8}
\renewcommand{\textfraction}{0.08}
\renewcommand{\floatpagefraction}{0.75}
\renewcommand{\contentsname}{목차}
\renewcommand{\abstractname}{초록}
\renewcommand{\refname}{참고문헌}
\renewcommand{\figurename}{그림}
\renewcommand{\tablename}{표}
\setcounter{tocdepth}{2}
\pagestyle{fancy}
\fancyhf{}
\fancyhead[L]{\small LATENT-ART-BENCH · 한국어 번역}
\fancyhead[R]{\small\texttt{f9e39354e462}}
\fancyfoot[C]{\thepage}
\renewcommand{\headrulewidth}{0.3pt}
\newcommand{\E}{\widehat{\mathcal E}}
\newcommand{\R}{\mathbb R}
\newcommand{\sourcefigure}[1]{%
  \IfFileExists{figures/#1}{%
    \includegraphics[width=\linewidth,height=0.72\textheight,keepaspectratio]{figures/#1}%
  }{%
    \fbox{\begin{minipage}{0.94\linewidth}
      \small\vspace{5pt}
      이 번역본에서는 그림 이미지 대신 원본 링크를 제공합니다.\par
      \href{https://github.com/isingmodel/latent-art-bench/blob/f9e39354e4625e070c100a9ef2d55c77552e2076/paper/figures/#1}{원본 그림 PDF 열기: \nolinkurl{#1}}
      \vspace{5pt}
    \end{minipage}}%
  }%
}
\title{텍스트-이미지 생성에서\\화가 이름이 유발하는 반응:\\공통 회화 효과를 넘어서}
\author{익명 저자}
\date{한국어 번역본 · 2026년 9월 15일}

\begin{document}
\maketitle
\thispagestyle{plain}
\begin{center}
\textit{Artist-Name Responses beyond a Shared Painting Effect\\in Text-to-Image Generation}
\end{center}
\begin{abstract}
화가의 이름은 화가 간 차이를 재현하지 않고도 이미지를 바꿀 수 있다. 여섯 생성 모델에 고정된 장면 14개를 제시하고, 화가 미지정, 일반 유화, 모네·시슬레·피사로·세잔의 화풍을 지정한 유화 프롬프트로 각각 두 번씩 생성하여 이미지 1,008개를 얻었다. 생성 이미지와 디지털 복제 이미지 649개에서 색채·공간·텍스처에 관한 31개 측정값을 비교했다. 각 화가를 네 화가의 평균과 비교함으로써 공통적인 변화와 화가 간 차이를 분리했다. 장면에 걸쳐 평균을 내고 반복 생성 잡음을 교정하면, 화가 미지정 프롬프트에서 화가 지정 프롬프트로 바뀔 때 발생하는 제곱 변화량의 82.5--95.7\%가 이름 간에 공통적이었다. 여기에는 유화 지시문의 기여도 포함된다. 배율 보정 전에는 모든 모델이 참조 자료의 차이와 정렬되지만, 화가를 전혀 구분하지 않는 예측기보다 우수하다고 입증된 모델은 없었다. 사후 진단에서는 다섯 모델의 모네--시슬레 정렬이 0에 가깝거나 반대 방향으로 나타났다. 배율 보정은 다른 장면에서 적합한 계수를 사용해 측정된 화가 대조값의 크기를 조정한다. 보정 전 오차 추정치는 FLUX.2 Max가 가장 낮았고, 홀드아웃 장면에서 평가한 보정 후 오차 추정치는 GPT Image 2가 가장 낮았다. 이미지를 변경한 것은 아니다. 이러한 측정만으로 지각적 충실도가 입증되지는 않는다.
\end{abstract}
\clearpage
{\small\linespread{0.95}\selectfont
\setlength{\parskip}{0pt}
\tableofcontents}
\clearpage

\section{서론}
‘모네의 화풍으로’ 그리라는 지시는 공통적인 회화적 외관을 도입하거나, 익숙한 색채·텍스처 단서를 과장하거나, 화가를 구별하는 특징을 바꿀 수 있다. 이러한 효과는 쉽게 혼동된다. 예를 들어 모네 작품 집합에 더 가까워졌다는 사실만으로 생성 모델이 모네와 시슬레의 차이를 재현한다고 볼 수는 없다. 일반적인 회화 지시만으로도 두 작품 집합 모두에 대한 근접도가 높아질 수 있기 때문이다.

본 연구는 화가 이름이 이름들 사이에 공통적으로 나타나는 반응을 넘어, \emph{출처가 기록된 디지털 복제 이미지}에서 관찰되는 화가 간 차이와 정렬되는 차이를 유발하는지 묻는다. 각 화가의 색채·공간·텍스처 평균 측정값을 네 화가의 평균과 비교하고, 생성 이미지가 이렇게 화가 라벨에 대응하는 차이를 재현하는지 평가한다. 평가 대상은 해당 참조 집합에 속하는 차이 패턴이다. 역사적 작품의 소재 구성과 복제 이미지 제작 과정도 화가 간 차이에 기여할 수 있다. 생성 장면의 텍스트를 고정하면 프롬프트 개입은 통제할 수 있지만, 역사적 작품들의 소재까지 동일해지는 것은 아니다.

이 구분은 평가에서 중요하다. 생성 모델은 화가 이름별 조건을 쉽게 분리할 수 있는 이미지를 만들면서도, 각 라벨에 해당하지 않는 참조 화가의 특징을 대응시킬 수 있다. 또한 참조 패턴을 따르는 변화는 맞추면서, 그 패턴을 따르지 않는 큰 차이를 추가할 수도 있다. 반대로 반응이 약한 모델이 강한 모델보다 제곱오차는 낮을 수 있지만, 모든 측면에서 화가를 더 정확히 구별하는 것은 아닐 수 있다. 단일 근접도 점수나 인식 정확도로는 이러한 가능성을 분리할 수 없다.

여섯 모델 구성에 동일한 장면과 프롬프트 지시문을 제공했다. 명시한 독립성 가정하에서 반복 생성 결과를 활용하여, 무작위 출력 변동을 체계적 오차로 계산하지 않으면서 불일치를 추정했다. 별도의 사후 점검에서는 화가 쌍의 차이, 대조값의 크기, 장면 의존성, 참조 자료의 품질을 살펴보았다.

\section{관련 연구}
\citet{somepalli2024style}은 학습 기반 화풍 기술자를 개발하고 확산 모델 간 유사성을 연구했다. \citet{su2025prompted}은 내용을 통제한 이름 치환을 사용해 프롬프트로 지정한 화가의 인식 성능을 평가하고, CLIP 공간에서 화가 지정·미지정 이미지를 실제 화가의 프로토타입과 비교했다. 본 연구의 질문은 공통 반응을 제거한 뒤 참조 화가들 사이에 남는 라벨별 차이에 관한 것이다. 이름 치환이나 모델 간 비교 자체가 새로운 것은 아니다.

\citet{deliege2025}은 전문가와 함께 화풍 모방과 정형화를 평가했다. AI-Pastiche는 진정성 판단과 프롬프트 준수를 연구하며\citep{asperti2025pastiche}, AI-WikiArt는 회화에서 도출한 캡션과 시각-언어 모델을 활용해 작가를 판별한다\citep{fu2025aiwikiart}. 이러한 과제는 설득력 있는 모방, 알아볼 수 있는 작가성, 분포의 일치를 구분한다. 본 연구의 주된 평가 대상은 측정된 색채·공간·텍스처 좌표에서 화가 라벨별 평균 대조값이 얼마나 일치하는지이다.

정량적 미술 연구는 색채, 구도, 이미지 통계를 다룬다\citep{kim2014,sigaki2018,lee2020}. \citet{kim2026}은 미술사적 분석과 맥락을 반영한 이미지-이미지 실험을 위해 조형적 표현과 맥락적 표현을 결합한다. 본 연구의 텍스트 전용 이름 개입에는 원작의 구도가 제공되지 않는다.

생성 모델 평가는 충실도와 다양성을 구분하며\citep{naeem2020}, 조건부 지표는 조건 내 변동과 평균 간 관계를 구별한다\citep{benny2021conditional}. 본 연구는 독립 분할 간 내적을 잡음 교정 제곱차의 주 추정량으로 사용하고\citep{diedrichsen2017representational}, 에너지 통계를 보완적인 분포 지표로 사용한다\citep{szekely2013}. 처리 방식에 대한 민감도\citep{parmar2022}와 학습 임베딩에서 분리도를 해석할 때의 한계\citep{asperti2026}를 고려해, 출처 품질을 점검하고 평가 대상을 명시된 디지털 참조 자료로 한정했다.

\section{실험 설계와 측정}\label{sec:design}
\subsection{화가와 참조 이미지}
주 분석 패널은 서로 관련된 네 화가의 서로 다른 복제 이미지 649개로 구성된다. 모네 297개, 시슬레 106개, 피사로 141개, 세잔 105개이다. 화가 평균에는 동일한 가중치를 부여하므로, 모네 작품 수가 더 많다고 해서 화가 간 비교의 평가 대상을 지배하지 않는다. 별도의 개발용 작품 221개로 공통 특징 단위를 설정했다. 이 변환을 적합하는 데 주 참조 이미지나 생성 이미지는 사용하지 않았다. 두 패널 모두 이번 실험 이전에 구축되어 선행 분석에 사용된 바 있다. 각 화가의 전체 작품에서 새롭게 확보해 숨겨 둔 표본이 아니라, 범위가 명시된 유한한 비교 패널이다.

참조 파일에는 Wikimedia Commons 출처 메타데이터와 Wikidata 작품 식별 정보가 포함되어 있다\citep{commonsglam,vrandecic2014wikidata}. 이러한 정보는 선정 과정과 권리 기록의 추적을 돕지만, 촬영·색 보정·보존 상태를 표준화하지는 않는다. 본 연구는 디지털 복제 이미지를 측정하므로, 이미지 획득 조건이 비교에 영향을 줄 수 있다.

\subsection{공통 장면을 사용하는 여섯 모델 실험}
GPT Image 1, GPT Image 2, GPT Image 2.5 Flare(이하 Flare), GPT Image 2.5 Sunburst(이하 Sunburst), Nano Banana 2, FLUX.2 Max(이하 FLUX)를 비교한다. 고정된 야외 장면 설명 14개는 수면, 건축 환경, 육지, 혼합 구도를 포괄한다. 각 장면에 여섯 지시문 조건을 교차 적용한다. 회화 지시 없음, ‘유화로 표현하라’(Render as an oil painting), 그리고 네 화가 각각에 대해 ‘[화가]의 화풍을 따른 유화로 표현하라’(Render as an oil painting in the style of [painter])이다. 이 외의 장면 텍스트는 동일하다. 모델--장면--지시문 조합마다 별도로 두 번 출력을 요청하여, 2026년 9월 10일 이미지 1,008개를 수집했다.

모든 모델에는 텍스트만 입력했으며, 출력은 $1{,}024\times1{,}024$픽셀 이미지이다. 네 OpenAI 모델에는 중간 품질 설정을 사용했고, 나머지 두 모델에는 각각 명시된 1K 구성을 사용했다. 이 비교는 기록된 구성에 관한 것이며, 각 모델의 최대 가능 품질이나 이미지당 동일 비용 조건을 비교하는 것은 아니다. 실제 서비스된 모델 버전은 독립적으로 검증하지 않았다(부록~\ref{app:provenance}). 정확한 프롬프트, 매개변수, 무작위 요청 순서는 수집 전에 고정했다. 장면 14개는 모든 가능한 내용에서 무작위로 추출한 표본이 아니라 고정 설계이다. 이전에 수행한 별도의 프롬프트 개입은 부록~\ref{app:cohorts}에 요약한다.

그림~\ref{fig:original-generated}은 한 장면에 대한 여섯 조건을 모두 보여준다. GPT Image 2의 세잔 조건에서 나무는 모네·시슬레 조건보다 더 짙은 청록색 윤곽을 보인다. 이는 이 장면에서 이름에 대한 반응이 눈에 보이는 사례이지만, 역사적 화가 간 차이와 일치하는지는 별도로 측정해야 한다.

\begin{figure}[H]
\centering\sourcefigure{specificity_original_generated.pdf}
\caption{하나의 장면, 여섯 프롬프트 조건, 여섯 생성 모델. 요청 장면은 자갈 강둑을 따라 굽이치는 넓은 강, 맞은편 강둑의 나무 세 그루, 흐린 하늘 아래 낮은 언덕이다. 각 모델 행에서 화가를 지정한 네 출력을 오른쪽의 화가 미지정·일반 유화 대조 조건과 비교한 뒤, 화가 이름 조건끼리 비교한다. 맨 위 행에는 표시를 위해 감사된 참조 영역을 제시하지만, 주 점수는 전체 프레임으로 계산한다. 이 작품들은 생성 모델의 입력으로 사용되지 않았으며 구도를 맞춘 작품도 아니다. 모든 생성 예시는 첫 번째 장면의 첫 번째 반복 출력이다. 전체 프롬프트, 선정 규칙, 출처는 부록~\ref{app:images}에 제시한다.}
\label{fig:original-generated}
\end{figure}

\subsection{특징이 측정하는 것}\label{sec:features}
표~\ref{tab:features}의 특징 31개는 반드시 함께 변하지는 않는 속성들을 포착한다. 두 이미지의 채도, 즉 색의 강도가 비슷해도 색상 다양성은 다를 수 있고, 전체 대비가 비슷해도 에지 방향과 미세한 텍스처는 다를 수 있다. 이러한 구분 덕분에 특징군별 결과를 해석할 수 있다. 다만 좌표 간 상관관계와 이미지 처리에 대한 민감도 때문에 유클리드 거리를 검증된 예술적 화풍 지표로 취급할 수는 없다.

\begin{table}[H]
\centering\small
\begin{tabularx}{\linewidth}{@{}>{\raggedright\arraybackslash}p{.15\linewidth}>{\raggedright\arraybackslash}p{.41\linewidth}>{\raggedright\arraybackslash}X@{}}
\toprule
특징군 & 측정 항목 & 해석\\\midrule
색채 (11) & 명도·채도의 중앙값과 IQR, 유채색 픽셀 비율, 색상 집중도·엔트로피, 세 간격에서의 색차와 그 기울기 & 밝기, 색의 강도와 다양성, 공간에 따라 색이 얼마나 빠르게 변하는지.\\[5pt]
공간 (8) & 스펙트럼 기울기·잔차·이방성, 방향 엔트로피, 수평--수직 균형, 그래디언트 중앙값·IQR, 사분면 발산도 & 거친 구조와 미세 구조의 구성, 지배적 방향, 에지 강도, 방향 분포의 영역별 변동.\\[5pt]
텍스처 (12) & 네 웨이블릿 에너지와 기울기·곡률, 세 국소 이진 패턴 엔트로피, 세 국소 변동계수 & 여러 스케일의 세부 정보, 국소 패턴의 다양성, 국소 휘도 대비.\\
\bottomrule
\end{tabularx}
\caption{고정된 디지털 특징 31개. 개수는 스칼라 좌표의 수이다. 텍스처 측정값은 물감의 물리적 요철을 측정하거나 붓질의 진위를 판별하는 값이 아니다. 정확한 스케일과 정의는 부록~\ref{app:features}에 제시한다.}
\label{tab:features}
\end{table}

이미지에는 방향 보정과, 유효한 내장 프로파일을 이용한 sRGB 변환을 적용한다. 프로파일이 없지만 호환되는 파일은 sRGB로 간주한다. 주 파이프라인은 종횡비를 유지하며 짧은 변을 512픽셀로 맞추고, 업샘플링 없이 Lanczos 리사이징을 사용한다. 회화 영역 마스크를 적용하지 않으므로 출처 이미지의 여백이나 색 보정 띠가 측정값에 포함될 수 있다. 각 좌표는 개발 패널의 중앙값과 사분위 범위(IQR)로 중심화·스케일링하며, 화가별 총가중치는 동일하게 한다. 따라서 한 단위는 해당 좌표의 개발 패널 IQR이지, 지각적 최소 감지 차이가 아니다. 유클리드 지표는 표준화 좌표에 동일한 가중치를 부여하며, 특징군별 가중치를 추가로 균등화하지 않는다.

\FloatBarrier
\section{화가 간 차이 비교}\label{sec:method}
\subsection{무엇을 일치로 보는가?}\label{sec:contrast-intuition}
비교는 세 단계로 이루어진다. 모든 이미지에서 동일한 속성 31개를 측정한다. 각 역사적 화가에 대해 네 화가의 전체 평균을 빼서 해당 화가가 집단과 어떻게 다른지 기록한다. 각 생성 장면의 화가 지정 출력 네 개에도 같은 평균 차감을 적용한 뒤, 대응하는 화가별 차이를 비교한다. 본 연구에서는 이 차이를 \emph{대조값(contrast)}이라고 부른다.

명도 중앙값을 예로 들면, 참조 자료에서 모네의 평균은 시슬레보다 개발 패널 IQR 단위로 $.325$ 낮다. GPT Image 2에서는 장면 14개와 두 반복에 걸친 차이가 같은 방향의 $.254$로 더 작다. 각 집합에서 네 화가의 평균을 빼더라도 이러한 간격은 유지된다. 공통적인 밝기 이동만으로는 이 구별을 만들 수 없다. 전체 비교에는 좌표 31개와 화가 네 명을 모두 사용한다.

두 점수는 서로 다른 질문에 답한다. \emph{정렬 반응} $\beta$는 생성 대조값이 참조 대조값 방향을 따라 얼마나 뻗어 있는지 측정하며, $\beta=1$이면 해당 성분의 강도가 일치한다. \emph{오차} $D$는 다른 방향의 차이까지 포함한 전체 불일치를 계산한다. 이상적인 생성 모델 세 개를 생각해 보자. 하나는 화가를 구분하지 않고, 하나는 참조 대조값을 정확히 일치시키며, 하나는 그 대조값을 두 배로 만든다. 이들의 기대 오차는 각각 1, 0, 1이다. 따라서 반응이 강할수록 오차가 커질 수도 있으며, 정렬 반응이 1에 가까워도 다른 부분에는 상당한 불일치가 남을 수 있다.

\subsection{두 점수의 계산}
$\mu_a\in\R^{31}$를 화가 $a$의 표준화된 참조 평균, $r_a=\mu_a-\frac14\sum_b\mu_b$를 그 대조값으로 둔다. 이 대조값들의 전체 제곱 크기 $H=\sum_a\|r_a\|^2$로 두 점수를 정규화한다. 여기서 $\|u\|^2$는 좌표 제곱의 합이며, $\langle u,v\rangle$는 대응 좌표 곱의 합이고, 모자 기호는 추정값을 나타낸다. 생성 측정값 $z_{msak}$의 첨자는 모델 $m$, 장면 $s$, 화가 $a$, 반복 $k\in\{1,2\}$를 나타낸다.
\begin{equation}
 d_{msak}=z_{msak}-\frac14\sum_b z_{msbk}.
 \label{eq:center}
\end{equation}
중심화는 공통적인 가산 이동을 제거하지만, 공통적인 배율 변화나 좌표 혼합은 여전히 대조값을 바꿀 수 있다(부록~\ref{app:score-details}). 장면 수 $S=14$일 때 정렬 반응은 다음과 같다.
\begin{equation}
 \widehat\beta_{msk}=\frac{\sum_a\langle d_{msak},r_a\rangle}{H},\qquad
 \widehat\beta_m=\frac1S\sum_s\frac{\widehat\beta_{ms1}+\widehat\beta_{ms2}}2.
 \label{eq:beta}
\end{equation}
양수는 전체적으로 참조 차이를 따르는 반응이 있음을 뜻하지만, 모든 화가 쌍에서 일치해야 한다는 뜻은 아니다. 1보다 크면 해당 방향의 반응이 과장되어 있음을 나타낸다. 오차를 측정할 때는 어느 한 반복의 불일치를 제곱하는 대신 두 반복의 불일치를 서로 곱한다.
\begin{equation}
 \widehat D_{ms}=\frac1H\sum_a
 \langle d_{msa1}-r_a,\ d_{msa2}-r_a\rangle,\qquad
 \widehat D_m=\frac1S\sum_s\widehat D_{ms}.
 \label{eq:distortion}
\end{equation}
조건부 평균이 안정적이고 반복 오차가 독립적이며 평균이 0이라면, 요청 간에 되풀이되지 않는 잡음의 기대 기여도는 0이다. 따라서 이 점수의 기댓값은 장면 조건부 평균 대조값의 정규화된 제곱오차가 된다. 개별 추정값은 음수일 수 있으며, 이를 그대로 유지하되 완벽한 복원보다 더 낫다고 해석하지 않는다. 대조값이 없는 $D=1$ 기준은 이론적인 예측기이지, 실제로 관찰한 일반 유화 조건이 아니다.

평가 대상은 모든 장면에서 동일한, 역사적 작품을 통합한 참조 대조값이다. 따라서 화가 간 차이가 장면마다 달라지는 모델은 장면 평균이 참조 자료와 일치해도 오차가 생길 수 있다. 그러한 변동이 반드시 예술적 오류인 것은 아니다. 본 연구는 장면 평균 후의 불일치와 추가되는 장면 변동 성분을 별도로 보고하며, 그 합이 $D$이다. 정확한 항등식과 유도는 부록~\ref{app:score-details}에 제시한다.

\subsection{모델 비교와 불확실성}
장면별로 모델을 비교한 뒤 대응된 오차 차이를 평균한다. 이렇게 공통 장면 설명을 통제하며, 특징·화가·개별 이미지를 독립 반복으로 취급하지 않는다.

새로운 특징 결과를 확인하기 전에 정렬 반응 여섯 개와 모델 쌍의 오차 차이 15개를 선언했다. 대응 장면별 Student 구간에는 이 21개 비교를 하나의 집합으로 묶어 Bonferroni 보정을 적용했으며, 근사적으로 95\% 동시 포함률을 갖는다. 구간이 0을 제외할 때 해당 비교의 차이가 구별된다. 절대 $D$ 구간은 명목 95\% 요약값으로, 다중 비교 보정을 하지 않은 기술적 구간이다.

이 구간들은 고정된 장면 설명 14개에 걸친 변동을 반영하며, 생성 잡음과 실제 장면 차이를 함께 포함한다. 이 특정 패널에서 새로 요청했을 때의 불확실성만 분리하지도 않고, 보지 못한 장면에 대한 포함률을 보장하지도 않는다. 출력 두 개로 잡음 교정은 가능하지만, 그 표본분포나 서비스의 독립성이 입증되는 것은 아니다.

사전 선언한 점검에서는 장면, 참조 표본, 특징군, 측정 창을 바꾼다. 별도로 선언한 민감도 분석에서는 제목에서 도출한 소재 범주 네 개의 비중을 각 화가 내에서 같게 맞춘다. 이러한 선택은 참조 패널과 지표에 대한 의존성을 점검하는 것이지, 역사적 작품의 구도를 일치시키는 것은 아니다. 전체 명세는 부록~\ref{app:estimation}에 남겨 둔다.

\subsection{추가 비교가 검증하는 것}\label{sec:diagnostics}
화가 미지정 조건과 일반 회화 조건은 중심화로 제거되는 공통 반응을 설명한다. 사후 진단에서는 어떤 화가 쌍이 전체 반응에 기여하는지, 대조값 크기가 오차를 어떻게 바꾸는지, 참조 자료 선택이 비교를 어떻게 바꾸는지라는 세 가지 질문을 다룬다. 크기 점검에서는 모델마다 하나의 음이 아닌 배율을 선택한다. 이 배율은 모든 화가 대조값과 좌표에 공통으로 적용되며, 장면 13개에서 오차가 최소화되도록 정한다. 제외한 한 장면에서 이를 평가하고, 장면 14개를 번갈아 제외하며 반복한다. 이 \emph{배율 보정}은 중심화된 생성 대조값에만 배율을 곱한다. 참조 대조값, 좌표 단위, 이미지는 그대로 유지한다. 이러한 진단은 사전 선언 결과를 해석하기 위한 것으로, 유의성 검정 집합을 확장하지 않는다. 수식은 부록~\ref{app:score-details}에, 대조 분석·적합·출처 감사는 부록~\ref{app:diagnostics}에 제시한다.

\section{결과}\label{sec:results}
\subsection{올바른 방향의 반응도 목표에서 벗어날 수 있다}
출력 1,008개 모두에서 유효한 측정값을 얻었다. 전체 프레임 참조 대상을 사용할 때 모든 모델의 정렬 반응 구간이 0보다 높았다(표~\ref{tab:specificity-models}). 따라서 명시한 추론 모형하에서는 공통적인 가산 반응만으로 화가 지정 조건을 설명하기에 충분하지 않다. 그러나 양의 정렬이 낮은 오차를 보장하지는 않는다. GPT Image 2의 참조 정렬 강도는 $\beta=.999\;[.893,1.104]$로 참조와 거의 일치하지만, 전체 오차는 $D=1.226$이다. 즉, 정렬된 성분 외에 상당한 차이가 추가된다.

\begin{table}[H]
\centering\small\setlength{\tabcolsep}{4pt}
\begin{tabularx}{\linewidth}{@{}>{\raggedright\arraybackslash}p{.24\linewidth}>{\raggedright\arraybackslash}Xr>{\raggedright\arraybackslash}X@{}}
\toprule
모델 & 정렬 반응 $\beta$ & \shortstack{보정 전\\오차 $D$} & FLUX 대비 $\Delta D$\\\midrule
GPT Image 1 & $0.938\;[0.709,1.168]$ & $1.572$ & $0.771\;[-0.037,1.580]$\\[3pt]
GPT Image 2 & $0.999\;[0.893,1.104]$ & $1.226$ & $0.425\;[-0.161,1.011]$\\[3pt]
GPT Image 2.5 Flare & $0.772\;[0.680,0.864]$ & $1.738$ & $0.937\;[0.256,1.618]$\\[3pt]
GPT Image 2.5 Sunburst & $0.824\;[0.680,0.968]$ & $1.641$ & $0.840\;[0.058,1.623]$\\[3pt]
Nano Banana 2 & $0.442\;[0.256,0.629]$ & $1.109$ & $0.309\;[-0.847,1.464]$\\[3pt]
FLUX.2 Max & $0.470\;[0.239,0.701]$ & $0.801$ & $0$ (기준)\\[3pt]
\bottomrule
\end{tabularx}
\caption{전체 프레임 참조 대상과의 일치. $\beta=1$이면 정렬된 진폭이 일치하며, 더 큰 값이 반드시 더 좋은 것은 아니다. $D$의 기댓값은 정확히 일치할 때 0, 화가를 구분하지 않을 때 1이다. 양의 $\Delta D$는 FLUX.2 Max에 유리하다. 대괄호는 사전 명시한 정렬 반응 여섯 개와 장면별 대응 모델 비교 15개의 집합에서 산출한 근사 95\% 동시 구간이다. 가장 낮은 FLUX의 $D$ 추정치에 대한 보정하지 않은 95\% 구간은 $[0.586,1.016]$으로, 화가 미구분 값 1을 포함한다. 이 구간은 기술적 요약이다.}
\label{tab:specificity-models}
\end{table}

FLUX.2 Max의 정렬 반응은 $.470$으로 더 작지만 오차 추정치는 $.801$로 가장 낮다. 보정된 비교에서는 Flare와 Sunburst의 오차가 FLUX보다 높다는 차이가 구별되며, 나머지 13개 구간은 0을 포함한다(부록~\ref{app:results}, 표~\ref{tab:all-specificity-pairs}). 이는 전체 프레임 참조를 사용한 결과이다. Sunburst에 대한 결과는 출처 교정 후 달라진다(\ref{sec:source-correction}절). 대조값이 없는 기준값 1보다 우수하다고 입증된 모델은 없다. 따라서 본 연구는 전체 불일치가 낮다는 사실보다 참조와 정렬되는 반응의 존재를 더 분명하게 검출한다.

그림~\ref{fig:specificity-geometry}에서 세잔은 시각적으로 분리되어 있는 반면, 생성된 모네·시슬레 평균은 서로 가까운 경우가 많다. M/S/P 확대 상자는 가까이 위치한 점들을 확대해 보여준다. 축은 표준화 특징 31개를 결합하며, 그 방향은 네 참조 평균만으로 적합한다. 두 축은 이 평균들 사이 변동의 86.99\%를 함께 설명한다. 아래의 쌍별 분석은 특징 31개 전체를 점검한다.

\begin{figure}[H]
\centering\sourcefigure{specificity_geometry.pdf}
\caption{화가 라벨별 패턴의 일치. 십자 표시는 참조 대조값이며, 빈 원은 장면 14개와 두 반복에 걸친 생성 대조값의 평균이다. 동일 라벨을 연결하는 선이 짧을수록 투영된 공간에서의 불일치가 작다. 흐린 점은 중심화된 장면·반복 대조값이다. 윤곽선으로 표시한 영역은 인접한 모네(M), 시슬레(S), 피사로(P) 상세도에서 공통 배율로 확대하며, 세잔은 전체도에 남겨 둔다. 각 시각화는 모델 간에 공통 축을 사용한다. 이 두 투영 축에서 일치한다고 해서 특징 31개 모두에서 일치하는 것은 아니다.}
\label{fig:specificity-geometry}
\end{figure}
\FloatBarrier

\subsection{장면 평균 변화의 대부분은 화가 이름들에 공통적이다}\label{sec:shared-change}
장면에 걸쳐 평균을 내면, 화가 지정 조건에서 화가 미지정 조건을 뺀 이동은 네 이름의 평균 이동과 화가 고유의 이탈로 이루어진다. 평균 이동에는 유화 지시와 이름에 대한 모든 공통 반응이 포함되며, 반복 교정 제곱 변화량의 82.5--95.7\%를 차지한다(표~\ref{tab:shared-change}). 평균을 내면 일관된 이동이 강조된다. 장면 내에서 분리하면 장면별 변동이 유지되어 공통 비율은 88.6--97.2\%가 된다(부록~\ref{app:estimation}). $\beta$는 참조 대조값 방향으로의 이탈을 측정한다.

\begin{table}[H]
\centering\small\setlength{\tabcolsep}{3pt}
\begin{tabularx}{\linewidth}{@{}l*{6}{>{\centering\arraybackslash}X}@{}}
\toprule
 & \shortstack{GPT Image\\1} & \shortstack{GPT Image\\2} & Flare & Sunburst & \shortstack{Nano\\Banana 2} & \shortstack{FLUX.2\\Max}\\\midrule
공통 성분 (\%) & 82.5 & 85.8 & 84.8 & 83.9 & 88.1 & 95.7\\
이름 간 성분 (\%) & 17.5 & 14.2 & 15.2 & 16.1 & 11.9 & 4.3\\\midrule
교정 코사인 & 0.689 & 0.774 & 0.766 & 0.837 & 0.930 & 0.916\\
제곱 크기 비 & 1.986 & 1.960 & 2.816 & 3.356 & 3.692 & 4.138\\
\bottomrule
\end{tabularx}
\caption{화가 미지정에서 화가 지정 프롬프트로 바뀔 때의 장면 평균·반복 교정 제곱 변화량 비율. 회화 지시문도 포함된다. 공통 성분은 네 이름의 평균 이동이고, 이름 간 성분은 그 평균으로부터의 이탈이다. 아래 두 행은 동일한 화가 미지정 기준을 사용해 공통 화가 지정 이동과 일반 유화 이동을 사후 비교한 값이다. 잡음 교정 코사인은 방향의 일치도를 측정하며 1은 평행을 뜻한다. 크기 비는 반복 교정 제곱 크기에서 공통 성분을 일반 유화 성분으로 나눈 값이다.}
\label{tab:shared-change}
\end{table}

공통 화가 지정 이동은 대체로 일반 유화 이동과 같은 방향을 가리키지만(잡음 교정 코사인 $.689$--$.930$; 1은 평행), 반복 교정 제곱 크기는 1.96--4.14배이다. 두 값 모두 장면 평균과 동일한 화가 미지정 기준을 사용하며, 모델별 값은 표~\ref{tab:shared-change}에 제시한다. 따라서 화가 이름들 사이에 공통적인 반응이 일반 회화 지시에 대한 반응과 같을 필요는 없다. 일반 유화 이동과 이름 지정으로 추가되는 공통 이동은 서로 강화하거나 상쇄할 수 있으므로, 두 제곱 크기는 독립적인 기여처럼 더할 수 없다(부록~\ref{app:estimation}).
\FloatBarrier

\subsection{전체 반응은 고르지 않은 화가 구별을 가린다}
사후 쌍별 진단은 \ref{sec:contrast-intuition}절의 명도 예시를 특징 31개 전체로 확장한다(그림~\ref{fig:artist-pairs}). 모네--시슬레 정렬 반응은 다섯 모델에서 $-.249$에서 $.129$ 사이인 반면, GPT Image 2에서는 $.402$이다. 그러나 이 쌍의 오차 추정치는 GPT Image 1이 더 낮다($.56$ 대 $1.16$). 더 강한 정렬이 더 낮은 불일치를 보장하지 않는 것이다. 세잔을 제외하면 FLUX의 전체 반응은 $.470$에서 $.096$으로, GPT Image 2의 반응은 $.999$에서 $.651$로 낮아진다.

모네와 시슬레를 서로 바꾸면 Flare, Sunburst, Nano Banana 2의 오차가 감소한다. 음의 방향으로 정렬된 쌍의 차이를 뒤집으면 크기를 바꾸지 않고 오차를 낮출 수 있다.

\begin{figure}[H]
\centering\sourcefigure{specificity_artist_pairs.pdf}
\caption{화가 쌍별 기술적 추정치이며, 쌍별 검정은 수행하지 않았다. M, S, P, C는 모네, 시슬레, 피사로, 세잔을 뜻한다. 쌍의 정렬 반응이 0이면 정렬 성분이 없고, 1이면 참조 강도와 일치하며, 1보다 크면 과장되고, 음수면 반대 방향이다. 기대 오차는 정확히 일치하면 0, 구별하지 않으면 1이다. 각 쌍에는 특징 31개 전체를 사용하고, 해당 참조 쌍 거리의 제곱으로 정규화한다. 음의 오차는 반복 교정 때문에 발생할 수 있으며, 정확한 일치보다 더 나은 상태를 뜻하지 않는다.}
\label{fig:artist-pairs}
\end{figure}
\FloatBarrier

\subsection{장면 평균과 대조값의 크기는 순서를 바꾼다}
장면별 적합도, 장면 평균 후의 적합도, 화가 대조값의 홀드아웃 배율 조정 후 적합도를 비교한다(표~\ref{tab:review-calibration}).

\emph{장면 평균.} Nano Banana 2의 평균 화가 대조값은 FLUX보다 오차가 낮다($.723$ 대 $.761$). 그러나 장면별로 평가하면 Nano Banana 2에 더 큰 변동 성분($.387$ 대 $.040$)이 추가되어 순서가 역전된다. $D$는 통합된 역사적 참조 대상 주변의 변동에 불이익을 주며, 평균을 내면 장면에 의존하는 차이가 상쇄될 수 있다.

\begin{table}[H]
\centering\small
\begin{tabularx}{\linewidth}{@{}l*{3}{>{\centering\arraybackslash}X}@{\hspace{1em}}>{\centering\arraybackslash}X@{}}
\toprule
모델 & \shortstack{보정 전\\장면 오차\\$D$} & \shortstack{장면 평균 후\\오차\\$D_{\rm aggregate}$} & \shortstack{장면\\변동\\$V_{\rm scene}$} & \shortstack{보정 후\\장면 오차\\$D_{\rm held}$}\\
\cmidrule(r){1-4}\cmidrule(l){5-5}
GPT Image 1 & 1.572 & 1.239 & 0.333 & 0.645\\
GPT Image 2 & 1.226 & 1.044 & 0.182 & \textbf{0.554}\\
2.5 Flare & 1.738 & 1.483 & 0.255 & 0.741\\
2.5 Sunburst & 1.641 & 1.337 & 0.305 & 0.706\\
Nano Banana 2 & 1.109 & \textbf{0.723} & 0.387 & 0.832\\
FLUX.2 Max & \textbf{0.801} & 0.761 & 0.040 & 0.714\\
\bottomrule
\end{tabularx}
\caption{전체 프레임 참조에 대한 사후 오차 분해와 배율 보정. $D=D_{\rm aggregate}+V_{\rm scene}$이다. 배율 보정은 나머지 장면 13개를 사용해 중심화된 생성 대조값의 크기를 조정한 뒤 제외한 장면에서 평가한다. 새로운 이미지를 생성하는 것은 아니다. 굵은 글씨는 각 오차 열의 최솟값이며, 유의성을 뜻하지 않는다.}
\label{tab:review-calibration}
\end{table}

\emph{대조값의 배율 조정.} GPT Image 2의 정렬 반응은 1에 가깝지만, 대조값의 제곱 크기는 참조의 두 배를 넘는다($Q=2.223$; 1이면 참조 크기와 일치한다. 부록~\ref{app:magnitude-coverage}). 축소하면 참조 정렬 성분은 약해지지만, 목표에서 벗어난 차이가 충분히 줄어 전체 오차 추정치가 낮아진다. GPT Image 2에 적합한 배율은 모든 대조 좌표에 약 0.45를 곱하며, 매번 다른 장면 13개만 사용해 정한다.

배율 보정 후에는 GPT Image 2의 오차 추정치가 $.554$로 가장 낮고, FLUX는 $.714$이다. 홀드아웃 폴드 14개 중 12개에서 GPT Image 2의 점수가 더 낮다. 어느 장면 하나를 제거하고 전체 절차를 다시 적합해도 평균적 우위는 유지된다(부록~\ref{app:magnitude-coverage}). 폴드들이 서로 겹치므로, 이 승리 횟수만으로 통계적 우월성이 입증되지는 않는다.

\subsection{참조 자료를 바꿔도 유지되는 비교는 무엇인가?}\label{sec:source-correction}
보정 전 장면 오차 $D$에서 FLUX는 장면 삭제, 측정 창, 참조 대상, 특징 가중치 민감도 분석 전반에 걸쳐 가장 낮은 점 추정치를 유지한다(부록~\ref{app:declared-sensitivities}, \ref{app:diagnostics}).

출처 처리 방식은 통계적 비교에 더 중요한 영향을 미친다. AI 보조 도구가 모든 참조·개발용 복제 이미지를 검사하여, 참조 이미지 90개와 개발 이미지 41개에서 주변부 인공물을 제외할 크롭 영역을 선정했다. 이 감사에는 사람의 검증을 사용하지 않았다. 크롭 후 개발용 스케일링을 다시 적합하면 FLUX의 오차 추정치는 $.801$에서 $.872$로 바뀐다. 교정 후에도 여섯 모델 모두 정렬 반응 구간이 0보다 높으며, GPT Image 2는 보정 후 가장 낮은 오차 추정치를 유지한다.

표~\ref{tab:source-comparison}는 달라진 비교를 명시적으로 보여준다. 교정된 영역과 다시 적합한 스케일링을 사용하면 Sunburst의 구간이 0을 가로지르고, Flare는 여전히 FLUX보다 오차가 높으며, GPT Image 1의 구간은 0보다 높은 쪽으로 이동한다. 마지막 결과는 사후 민감도 분석이지 새로운 확증적 발견이 아니다. 출처 교정 후에도 FLUX의 점 추정치는 가장 낮지만, 특정 모델 차이를 뒷받침하는 근거는 달라진다.

\begin{table}[H]
\centering\small
\begin{tabularx}{\linewidth}{@{}Xcc@{}}
\toprule
모델 비교 & \shortstack{전체 프레임 차이\\{[보정 구간]}} & \shortstack{크롭 + 스케일링 교정\\차이 [보정 구간]}\\\midrule
GPT Image 1 $-$ FLUX & $0.771\;[-0.037,1.580]$ & $0.758\;[0.026,1.491]$\\[2pt]
2.5 Flare $-$ FLUX & $0.937\;[0.256,1.618]$ & $0.864\;[0.193,1.535]$\\[2pt]
2.5 Sunburst $-$ FLUX & $0.840\;[0.058,1.623]$ & $0.705\;[-0.026,1.436]$\\[2pt]
\bottomrule
\end{tabularx}
\caption{감사된 회화 영역과 다시 적합한 개발용 스케일링을 사용한 출처 교정 전후의 배율 보정 전 오차 차이. 양의 차이는 FLUX.2 Max에 유리하다. 전체 프레임 구간은 평가 항목 21개에 대한 95\% 동시 절차를 사용한다. 교정 구간은 이 절차를 사후에 다시 적용한 것으로, 확증적 검정을 추가하지 않는다.}
\label{tab:source-comparison}
\end{table}

별도의 점검에서는 혼합 장면 세 개를 제외하고 수면·건축·육지 장면 11개에 범주별 참조 평균을 사용한다. 제목에서 도출한 참조 라벨을 시각 라벨로 대체해도 FLUX의 오차 추정치가 가장 낮다. 대조 분석과 감사 내용은 부록~\ref{app:diagnostics}에 제시한다.

\section{논의}\label{sec:discussion}
이름에 대한 반응을 발견하는 것과 화가 간 차이를 일치시키는 것은 별개의 성과이다. FLUX의 보정 전 오차 추정치가 가장 낮다는 사실과 모네--시슬레 정렬이 약하다는 사실은 동시에 성립한다. 따라서 모델 비교에서는 명시한 참조 평가 대상에 대해 전체 점수와 쌍별 점수를 함께 보고하고, 추정치와 통계적으로 구별되는 비교를 구분해야 한다. 또한 장면별 점수인지, 평균 패턴의 점수인지, 다른 장면에서 정한 배율을 홀드아웃 대조값에 적용한 점수인지 명시해야 한다. 이러한 선택에 따라 모델 순서가 달라지기 때문이다.

이 개입이 식별하는 것은 프롬프트 조건 간 차이이지, 그 차이의 내부 원인은 아니다. 학습 데이터 구성, 가이던스, 공통 변환 등은 여전히 가능한 설명으로 남는다.

\subsection{범위와 남은 한계}
평가 대상은 유한한 디지털 참조 집합의 차이 패턴이다. 여기에는 작품 자체의 속성뿐 아니라 소재 선택과 복제 이미지 제작 과정이 섞여 있다. 거친 내용 범주만으로는 구도, 시점, 작가의 활동 시기, 생성 이미지의 실제 장면 준수 여부를 일치시킬 수 없다. 출처 감사는 눈에 보이는 주변부 영역과 제목·시각 범주의 불일치를 다루지만, AI 판단, 불확실한 경계, 남아 있는 획득 조건 차이 때문에 교정에는 한계가 있다. 작품당 복제 이미지가 하나뿐이므로 조명, 샤프닝, 보존, 디지털화 출처의 효과를 분리할 수 없다.

좌표 31개는 디지털 색채·공간 구성·텍스처를 기술할 뿐, 물리적 붓질이나 지각되는 화가 유사성을 나타내지 않는다. 수치적 재실행과 좌표 가중치에 대한 강건성만으로 그러한 판단이 검증되지는 않는다. 추론은 또한 서로 관련된 화가 네 명, 작성된 장면 설명 14개, 두 반복, 하나의 지시문 템플릿에 조건부이다. 시뮬레이션은 명시한 잡음 모형을 점검하는 것이지, 실제 서비스의 독립성을 점검하는 것은 아니다. 이미지 패널은 비교를 직접 살펴볼 수 있게 하지만, 정확히 같은 픽셀에 대한 공개 접근이 불완전하므로 독립적인 측정 재현에는 여전히 제약이 있다.

\section{결론}
장면 평균 후에는 화가 미지정 프롬프트에서 화가 지정 프롬프트로의 특징 제곱 변화량 대부분이 이름 사이에 공통적이며, 여기에는 유화 지시의 기여가 포함된다. 그럼에도 여섯 모델 모두 참조 자료와 정렬되는 화가 간 차이를 보인다. 그러나 배율 보정 전에는 화가를 구분하지 않는 오차 기준보다 우수하다고 입증된 모델이 없다. 화가 쌍별 일치가 고르지 않고, 배율 보정 후 오차 순서가 바뀐다는 점은 단일 전체 점수만으로는 충분하지 않은 이유를 보여준다.

\section*{데이터 및 코드 공개}
\addcontentsline{toc}{section}{데이터 및 코드 공개}
\href{https://github.com/isingmodel/latent-art-bench}{저장소}에는 프로토콜, 프롬프트, 요청 설정, 출처 메타데이터, 스케일러, 특징 벡터가 포함되어 있다. 주 분석, 진단 확장, 출처 품질 민감도 분석은 각각 입력 해시와 재실행 코드를 갖춘 별도의 버전 기록으로 남아 있다. 출처 감사에는 모든 작품의 식별 정보, 영역 상자, 라벨 결정, 교정 벡터가 기록되어 있으며, 원래 결과는 그대로 유지된다.

그림~\ref{fig:original-generated}은 생성 예시 36개와 참조 예시 네 개를 보여준다. 부록~\ref{app:images}는 선정 방식을 설명하며, 함께 제공되는 매니페스트에는 식별 정보와 바이트 해시가 기록되어 있다. 공개된 수치 재실행은 벡터에서 시작한다. 독립적으로 특징을 다시 추출하려면 정확히 같은 출처 픽셀이 필요하지만, 이 픽셀들은 경량 저장소 외부에 보관되어 있다. 이전 실험에는 \href{https://github.com/isingmodel/latent-art-bench/releases}{공개 릴리스}가 있다. 현재 실험과 교정 결과에는 재배포가 허용되는 정확한 픽셀 또는 검증된 복구 경로를 포함한 전용 공개 아카이브가 아직 필요하다.

% 참고문헌을 파일 안에 포함하므로 BibTeX 실행이나 외부 .bib 파일이 필요하지 않습니다.
\begingroup
\small
\setlength{\bibsep}{5pt}
\begin{thebibliography}{20}
\addcontentsline{toc}{section}{참고문헌}
\bibitem[Asperti(2026)]{asperti2026}
Andrea Asperti.
On the separation of human and AI-generated images in CLIP embedding space.
arXiv:2608.25609, version 1, 2026.
URL \url{https://arxiv.org/abs/2608.25609}. Preprint.

\bibitem[Asperti et~al.(2025)]{asperti2025pastiche}
Andrea Asperti, Franky George, Tiberio Marras, Razvan Ciprian Stricescu, and Fabio Zanotti.
A critical assessment of modern generative models' ability to replicate artistic styles.
\textit{Big Data and Cognitive Computing}, 9(9):231, 2025.
doi: \href{https://doi.org/10.3390/bdcc9090231}{10.3390/bdcc9090231}.

\bibitem[Benny et~al.(2021)]{benny2021conditional}
Yaniv Benny, Tomer Galanti, Sagie Benaim, and Lior Wolf.
Evaluation metrics for conditional image generation.
\textit{International Journal of Computer Vision}, 129:1712--1731, 2021.
doi: \href{https://doi.org/10.1007/s11263-020-01424-w}{10.1007/s11263-020-01424-w}.
URL \url{https://link.springer.com/article/10.1007/s11263-020-01424-w}.

\bibitem[Deliège et~al.(2025)]{deliege2025}
Adrien Deliège, Jeanne Marlot, Marc Van Droogenbroeck, and Maria Giulia Dondero.
How good is the machine at the imitation game? On stylistic characteristics of AI-generated images.
\textit{Journal of Imaging}, 11(12):429, 2025.
doi: \href{https://doi.org/10.3390/jimaging11120429}{10.3390/jimaging11120429}.
URL \url{https://pmc.ncbi.nlm.nih.gov/articles/PMC12734345/}.

\bibitem[Diedrichsen and Kriegeskorte(2017)]{diedrichsen2017representational}
Jörn Diedrichsen and Nikolaus Kriegeskorte.
Representational models: A common framework for understanding encoding, pattern-component, and representational-similarity analysis.
\textit{PLOS Computational Biology}, 13(4):e1005508, 2017.
doi: \href{https://doi.org/10.1371/journal.pcbi.1005508}{10.1371/journal.pcbi.1005508}.

\bibitem[Fu et~al.(2025)]{fu2025aiwikiart}
Tarian Fu, Javier Conde, Gonzalo Martínez, Pedro Reviriego, Elena Merino-Gómez, and Fernando Moral.
Artificial intelligence and misinformation in art: Can vision language models judge the hand or the machine behind the canvas?
arXiv:2508.01408, version 1, 2025.
URL \url{https://arxiv.org/abs/2508.01408v1}.

\bibitem[Kim et~al.(2014)]{kim2014}
D. Kim, S.-W. Son, and H. Jeong.
Large-scale quantitative analysis of painting arts.
\textit{Scientific Reports}, 4:7370, 2014.
doi: \href{https://doi.org/10.1038/srep07370}{10.1038/srep07370}.

\bibitem[Kim et~al.(2026)]{kim2026}
J. Kim, B. Lee, T. You, and J. Yun.
Context-aware multimodal AI navigates hidden pathways in five centuries of art evolution.
\textit{Proceedings of the National Academy of Sciences}, 123(30):e2517969123, 2026.
doi: \href{https://doi.org/10.1073/pnas.2517969123}{10.1073/pnas.2517969123}.

\bibitem[Lee et~al.(2020)]{lee2020}
B. Lee, M. K. Seo, D. Kim, I.-S. Shin, M. Schich, H. Jeong, and S. K. Han.
Dissecting landscape art history with information theory.
\textit{Proceedings of the National Academy of Sciences}, 117(43):26580--26590, 2020.
doi: \href{https://doi.org/10.1073/pnas.2011927117}{10.1073/pnas.2011927117}.

\bibitem[Mallat(1989)]{mallat1989wavelet}
Stéphane G. Mallat.
A theory for multiresolution signal decomposition: The wavelet representation.
\textit{IEEE Transactions on Pattern Analysis and Machine Intelligence}, 11(7):674--693, 1989.
doi: \href{https://doi.org/10.1109/34.192463}{10.1109/34.192463}.
URL \url{https://www.di.ens.fr/~mallat/papiers/MallatTheory89.pdf}.

\bibitem[Naeem et~al.(2020)]{naeem2020}
Muhammad Ferjad Naeem, Seong Joon Oh, Youngjung Uh, Yunjey Choi, and Jaejun Yoo.
Reliable fidelity and diversity metrics for generative models.
In \textit{Proceedings of the 37th International Conference on Machine Learning}, volume 119 of \textit{Proceedings of Machine Learning Research}, pages 7176--7185, 2020.
URL \url{https://proceedings.mlr.press/v119/naeem20a.html}.

\bibitem[Ojala et~al.(2002)]{ojala2002lbp}
Timo Ojala, Matti Pietikäinen, and Topi Mäenpää.
Multiresolution gray-scale and rotation invariant texture classification with local binary patterns.
\textit{IEEE Transactions on Pattern Analysis and Machine Intelligence}, 24(7):971--987, 2002.
doi: \href{https://doi.org/10.1109/TPAMI.2002.1017623}{10.1109/TPAMI.2002.1017623}.

\bibitem[Parmar et~al.(2022)]{parmar2022}
Gaurav Parmar, Richard Zhang, and Jun-Yan Zhu.
On aliased resizing and surprising subtleties in GAN evaluation.
In \textit{IEEE/CVF Conference on Computer Vision and Pattern Recognition}, 2022.
URL \url{https://arxiv.org/abs/2104.11222}.

\bibitem[Sharma et~al.(2005)]{sharma2005ciede}
Gaurav Sharma, Wencheng Wu, and Edul N. Dalal.
The CIEDE2000 color-difference formula: Implementation notes, supplementary test data, and mathematical observations.
\textit{Color Research and Application}, 30(1):21--30, 2005.
doi: \href{https://doi.org/10.1002/col.20070}{10.1002/col.20070}.

\bibitem[Sigaki et~al.(2018)]{sigaki2018}
H. Y. D. Sigaki, M. Perc, and H. V. Ribeiro.
History of art paintings through the lens of entropy and complexity.
\textit{Proceedings of the National Academy of Sciences}, 115(37):E8585--E8594, 2018.
doi: \href{https://doi.org/10.1073/pnas.1800083115}{10.1073/pnas.1800083115}.

\bibitem[Somepalli et~al.(2024)]{somepalli2024style}
Gowthami Somepalli, Anubhav Gupta, Kamal Gupta, Shramay Palta, Micah Goldblum, Jonas Geiping, Abhinav Shrivastava, and Tom Goldstein.
Investigating style similarity in diffusion models.
In \textit{European Conference on Computer Vision}, 2024.
URL \url{https://www.ecva.net/papers/eccv_2024/papers_ECCV/papers/08294.pdf}.

\bibitem[Su et~al.(2025)]{su2025prompted}
Grace Su, Sheng-Yu Wang, Aaron Hertzmann, Eli Shechtman, Jun-Yan Zhu, and Richard Zhang.
Identifying prompted artist names from generated images, 2025.
URL \url{https://arxiv.org/abs/2507.18633v1}. Version 1.

\bibitem[Székely and Rizzo(2013)]{szekely2013}
Gábor J. Székely and Maria L. Rizzo.
Energy statistics: A class of statistics based on distances.
\textit{Journal of Statistical Planning and Inference}, 143(8):1249--1272, 2013.
doi: \href{https://doi.org/10.1016/j.jspi.2013.03.018}{10.1016/j.jspi.2013.03.018}.

\bibitem[Vrandečić and Krötzsch(2014)]{vrandecic2014wikidata}
Denny Vrandečić and Markus Krötzsch.
Wikidata: A free collaborative knowledgebase.
\textit{Communications of the ACM}, 57(10):78--85, 2014.
doi: \href{https://doi.org/10.1145/2629489}{10.1145/2629489}.

\bibitem[Wikimedia Commons(2026)]{commonsglam}
Wikimedia Commons.
Commons:GLAM, 2026.
URL \url{https://commons.wikimedia.org/wiki/Commons:GLAM}.
Accessed 11 September 2026.
\end{thebibliography}
\endgroup
\clearpage

\appendix
\section{특징의 정의}\label{app:features}
색채 좌표에는 D65 광원과 $2^\circ$ 표준 관찰자 조건의 CIELAB을 사용한다. 명도와 채도에서 각각 중앙값과 IQR을 하나씩 얻는다. 유채색 비율은 $C^*\geq5$인 픽셀의 비율이다. 색상 집중도는 채도로 가중한 원형 평균의 크기이며, 색상 엔트로피는 동일 간격의 구간 24개를 사용하고 구간 수의 로그로 정규화한다. 두 값 모두 채도 가중치를 사용하며, 유채색 픽셀이 1\% 미만이면 0으로 둔다. 세 간격, 즉 짧은 변 길이의 1\%, 4\%, 16\% 각각에서 수평·수직 CIEDE2000 차이를 합친 값의 중앙값을 구한다. 로그--로그 기울기로 스케일 의존성을 요약한다\citep{sharma2005ciede}.

공간·텍스처 좌표에는 선형 sRGB 휘도를 사용한다. 푸리에 분석 전 Tukey 창(매개변수 $.1$)을 적용한다. 짧은 변당 4--128주기를 로그 간격의 반지름 구간 32개로 나누고, Theil--Sen 직선으로 스펙트럼 기울기와 제곱평균제곱근 잔차를 구한다. 파워로 가중한 축방향 집중도로 이방성을 측정한다. Scharr 미분으로 그래디언트 중앙값과 IQR, 크기로 가중한 18구간 방향 엔트로피, 수평--수직 균형을 계산한다. 전체 히스토그램에 대한 각 사분면 방향 히스토그램의 Jensen--Shannon 발산도를 평균하여 영역별 변동을 측정한다.

4단계 비다운샘플링 \texttt{db2} 웨이블릿 변환에서 세부 계수의 평균제곱 에너지 네 개에 로그를 취하고, 이어서 그 선형 기울기와 평균 이차 차분을 계산한다\citep{mallat1989wavelet}. 균일 국소 이진 패턴 엔트로피는 양자화된 휘도에 $(P,R)=(8,1),(16,2),(32,4)$를 적용한다\citep{ojala2002lbp}. 세 국소 변동계수의 중앙값을 구할 때는 짧은 변의 1\%, 4\%, 16\%를 반올림해 창 너비로 사용하고, 짝수 너비는 다음 홀수로 늘린다. 평균의 하한은 $10^{-6}$으로 둔다. 고정된 구현에는 경계 처리와 안정화 상수가 기록되어 있다. 별도의 스케일러는 모네 101점, 시슬레 36점, 피사로 48점, 세잔 36점으로 구성된다. 각 가중 중앙값과 IQR을 계산할 때 화가별 총가중치를 같게 한다.

\section{추정량의 세부 사항과 민감도 분석}\label{app:estimation}
$\theta_{msa}=\mathbb E[d_{msak}]$로 두고 $d_{msak}=\theta_{msa}+\eta_{msak}$로 나타내자. 반복 오차의 평균이 0이고 $\mathbb E\langle\eta_{msa1},\eta_{msa2}\rangle=0$이면 다음이 성립한다.
\[
 \mathbb E[\widehat D_{ms}]
 =H^{-1}\sum_a\|\theta_{msa}-r_a\|^2.
\]
중심화는 한 반복 안에서 화가 간 의존성을 만들지만, 반복 블록이 독립적이면 반복 간 교차항은 여전히 사라진다. 공통 반응의 경우 $\theta_{msa}=0$이므로, 화가 지정 이미지 집합 전체를 참조 자료 쪽으로 이동시키더라도 기대 $D$는 1이다. 방향은 정확히 정렬되지만 배율이 달라진 반응 $\theta_{msa}=br_a$의 기대 오차는 $D=(b-1)^2$이다. 직교 방향의 이탈은 그 제곱 크기만큼 추가된다. 이 항등식은 화가 대조값이 더 강하다고 해서 참조 대비 제곱오차가 더 낮아지는 것은 아닌 이유를 설명한다.

통합 오차 $D_{\rm aggregate}$는 장면 평균 대조값에 교차곱 추정량을 적용한 것이다. 식~\ref{eq:scene-decomposition}는 이를 주 분석의 장면별 동일 가중 오차와 연결한다. 화가에도 동일한 가중치를 부여한다. 각 화가의 가산 기여도를 합하면 $D$가 되지만, 중심화가 네 조건을 결합하므로 이러한 기여도는 독립적인 화가 충실도 점수가 아니다.

모델 간 대조에서는 완전한 대응 장면마다 오차 차이를 하나씩 계산한다. 그 표준오차는 차이들의 표본 표준편차를 $\sqrt n$으로 나눈 값이다. 동시 구간의 반폭은 $t_{n-1,1-.05/(2\cdot21)}\,\mathrm{SE}$이다. 여섯 정렬 반응 구간에도 같은 비교 집합 보정을 사용한다. 추가적인 명목 구간과 대응 장면 부트스트랩 5,000회는 기술적 점검이지, 유의성에 도달하기 위한 대안 경로가 아니다. 참조 민감도 분석에서는 화가별로 독립적인 부트스트랩 표본 1,000개를 추출하고 참조 평균과 $H$를 다시 계산하되, 생성 벡터와 스케일링은 고정한다. 이는 모든 불확실성을 함께 다루는 것이 아니라 참조 패널 의존성을 기술한다.

새로운 특징 결과를 확인하기 전에 선언한 네 범주 참조 민감도 분석에서는 고정 어휘집을 출처 제목에 적용한다. 범주는 수면 중심, 건축 환경 중심, 경로 중심, 탁 트인 땅 또는 수목이 있는 육지이다. 각 화가 내에서 범주 평균마다 4분의 1의 가중치를 부여한다. 화가·범주 층 내 재표집은 이 라벨에 조건부이며, 분류 오차는 반영하지 않는다. 이전의 통제된 이름 지정 패널에는 아래에서 설명하는 별도의 세 범주 시각 코딩 체계를 사용했다.

공통 효과 진단에서는 각 반복 내에서 조건별로 장면 평균을 구한다. 화가 지정에서 미지정을 뺀 차이를 $h_{ak}$, $c_k=\frac14\sum_a h_{ak}$로 두면, 공통 성분과 고유 성분의 제곱 크기는 각각 $4\langle c_1,c_2\rangle$와 $\sum_a\langle h_{a1}-c_1,h_{a2}-c_2\rangle$이다. 전체가 양수일 때 공통 비율을 보고하며, 교정된 성분은 음수일 수 있다. 장면 내 버전에서는 각 장면에서 두 제곱 변화 성분을 따로 계산하고, 각 성분을 장면에 걸쳐 평균한 다음 그 비를 구한다. 이때 공통 비율은 88.6\%(GPT Image 1)에서 97.2\%(FLUX) 사이이다.

일반 유화에서 화가 미지정을 뺀 이동을 $g_k$라고 하자. 통상적인 반복 평균 코사인은 두 벡터에 잡음이 있는 동일한 미지정 추정값을 공유하므로, 분자의 기댓값에 그 오차 공분산의 대각합이 더해진다. 대신 \ref{sec:shared-change}절에 보고한 교정 코사인은 두 노름 제곱이 모두 양수일 때 반복 간 교차곱을 사용한다.
\[
 \frac{\langle g_1,c_2\rangle+\langle g_2,c_1\rangle}
 {2\sqrt{\langle g_1,g_2\rangle\langle c_1,c_2\rangle}}.
\]
이 비 자체는 불편 추정량이 아니며 $[-1,1]$로 제한되지도 않는다. 주 분석의 $D$는 화가 지정 조건만 사용하므로 공통으로 공유하는 화가 미지정 대조 조건의 잡음에 영향을 받지 않는다.

\subsection{추가 점수 항등식}\label{app:score-details}
공통 아핀 변환 $z\mapsto Az+b$하에서 중심화는 가산항 $b$를 제거하지만, 변환된 대조값 $Ad$는 남긴다.

추정량에는 정확한 분해식이 존재한다. $e_{msak}=d_{msak}-\widehat\beta_{msk}r_a$로 쓰면 다음과 같다.
\begin{equation}
 \widehat D_{ms}=
 (\widehat\beta_{ms1}-1)(\widehat\beta_{ms2}-1)
 +\frac1H\sum_a\langle e_{msa1},e_{msa2}\rangle.
 \label{eq:decomposition}
\end{equation}
첫 항은 참조 정렬 진폭 오차를 측정한다. 잔차는 참조 배열을 하나로 이어 붙인 전체 벡터에 직교한다. 이것이 반드시 참조 평균들이 특징 공간에서 생성하는 부분공간 밖에 있다는 뜻은 아니며, 진정한 예술적 변동 전체의 범위 밖에 있다는 뜻은 더더욱 아니다. 유한 표본에서는 두 항 모두 음수가 될 수 있다.

$\theta_s$를 네 조건부 평균 대조값을 이어 붙인 벡터, $\bar\theta=S^{-1}\sum_s\theta_s$라고 하자. 모집단 평가 대상은 다음과 같이 분해된다.
\begin{equation}
 D_{\rm scene}=\frac{\|\bar\theta-r\|^2}{H}
 +\frac1{SH}\sum_s\|\theta_s-\bar\theta\|^2
 =D_{\rm aggregate}+V_{\rm scene}.
 \label{eq:scene-decomposition}
\end{equation}
대조값의 강도와 방향을 구분하기 위해 다음을 정의한다.
\begin{equation}
 \widehat Q_m=\frac1{SH}\sum_{s,a}\langle d_{msa1},d_{msa2}\rangle,
 \qquad \widehat D_m=1-2\widehat\beta_m+\widehat Q_m.
 \label{eq:magnitude}
\end{equation}
$\widehat Q_m>0$일 때 $\widehat\beta_m/\sqrt{\widehat Q_m}$는 잡음 교정 정렬 요약값이며, 유한 표본 추정치는 $[-1,1]$ 안에 있을 필요가 없다. 중심화된 모든 생성 대조값에 $c$를 곱하면 $\widehat D_m(c)=1-2c\widehat\beta_m+c^2\widehat Q_m$가 된다. 나머지 장면 13개에서 $c=\max(0,\widehat\beta/\widehat Q)$를 적합하고 홀드아웃 장면에서 평가한다. 이 적합 기준은 학습에 사용한 $\widehat Q$가 양수일 때만 사용할 수 있으며, 관찰된 모든 폴드에서 그러했다. 이 반사실적 조정은 특징 벡터에 관한 것이지, 실제로 구현한 이미지 변환에 관한 것이 아니다.

\subsection{대응 비교와 사전 선언한 민감도 분석}\label{app:declared-sensitivities}
다음 장면 고정효과 회귀를 적합한다.
\begin{equation}
 \widehat D_{ms}=\alpha_s+\gamma_m+\varepsilon_{ms},
 \qquad\gamma_{\text{GPT Image 2}}=0.
 \label{eq:regression}
\end{equation}
균형 패널에서는 계수 차이가 추론에 사용한 대응 장면 평균과 같다. 회귀 기준을 바꿔도 이러한 차이는 변하지 않는다. 표~\ref{tab:specificity-models}에서는 FLUX를 기준으로 제시한다. 사전 선언한 민감도 분석은 각 장면을 하나씩 삭제하고, 장면과 참조 작품을 재표집하며, 특징군을 살펴보고 텍스처를 제외한다. 중앙 정사각형 창은 생성·참조 이미지의 주변부 내용을 버려 종횡비를 통제한다. 모든 장면 삭제에서 FLUX가 가장 낮은 보정 전 오차를 유지하며($D=.754$--$.852$), 정사각형 창에서도 그렇다($D=.765$; 다음으로 낮은 Nano Banana 2는 $1.024$).

\subsection{분포 근접도와 반복 산포}
화가 $a$에 대해 \emph{에너지 불일치}는 생성 집합 $Y_a$와 참조 집합 $X_a$를 다음과 같이 비교한다.
\begin{equation}
\begin{split}
 \E(X,Y)={}&\frac{2}{n_Xn_Y}\sum_{i,j}\|x_i-y_j\|\\
 &-\frac1{n_X^2}\sum_{i,i'}\|x_i-x_{i'}\|
 -\frac1{n_Y^2}\sum_{j,j'}\|y_j-y_{j'}\|.
\end{split}
 \label{eq:energy}
\end{equation}
유한 표본에서의 기댓값은 집합 내 산포에 의존하므로, 생성 표본 수가 같아도 비교 편향이 없어지지는 않는다. 반복 장면 설계에 대한 교정은 부록~\ref{app:diagnostics}에 제시한다. 장면 내 반복 산포 $\frac1{2S}\sum_s\|z_{msa1}-z_{msa2}\|^2$는 두 출력의 표본분산 대각합을 평균한 값이다. 이를 통해 반복 산포를 장면 설명 간 변동과 구분한다.

\section{사후 평가 대상 진단}\label{app:diagnostics}
이 사후 분석들은 수집이 완료된 생성 이미지 1,008개와 원래의 추론 비교 집합을 그대로 유지한다. 초기 진단에는 변경하지 않은 참조 벡터와 스케일링을 사용하며, 이후 출처 감사에서는 회화 영역, 개발용 스케일링, 내용 라벨을 별도로 바꾼다.

\subsection{내용별 평가 대상과 실제 회화 대조 분석}
기록된 참조 범주는 고정된 출처 제목 어휘집으로 부여한 수면 중심, 건축 환경 중심, 경로 중심, 탁 트인 땅 또는 수목이 있는 육지이다. 이 순서에 따른 작품 수는 모네 $(191,67,8,31)$, 시슬레 $(52,29,9,16)$, 피사로 $(28,77,12,24)$, 세잔 $(31,28,12,34)$이다. 범주별 평균도 통합 참조 대상과 마찬가지로 화가들에 걸쳐 중심화한다. 통합 대상에서 벗어나는 평균 제곱 크기는 $.438H$이며, 이는 유한 표본 오차를 포함한 관찰 참조 평균의 기술값이다.

생성 비교에서는 수면 장면 네 개, 건축 장면 네 개, 육지 장면 세 개에 대응하는 범주 평균을 사용하고, 혼합 장면 세 개는 제외한다. 이 $S=11$개 장면에서 범주별 교차곱 오차를 평균한 뒤 $H_c=S^{-1}\sum_s\sum_a\|r_{a,c(s)}\|^2=7.413$으로 나눈다. 통합 참조에서는 $H=5.915$이다. 따라서 표~\ref{tab:review-targets}의 각 열은 고유한 분모를 사용한다. 변경하지 않은 11개 장면의 통합 비교는 장면 부분집합 선택의 효과를 분리한다. 포함된 세 범주에서 제목 기반 층별 참조 작품 수는 화가·범주당 16--191점이지만, 시각적으로 장면이 일치하는 역사적 작품을 제공하는 것은 아니다.

실제 회화 대조 분석 1,000회 각각에서 화가·범주 셀의 작품들을 무작위로 겹치지 않는 두 절반으로 나누고, 더 작은 절반을 참조 대상 추정에 사용한다. 통합 대조 분석은 유사 장면 14개 각각에 대해 각 화가의 홀드아웃 절반에서 작품 두 점을 복원 방식으로 독립 추출한다. 범주 층화 대조 분석은 수면·건축·육지 유사 장면 11개 각각에 대해 홀드아웃 범주 안에서 표집한다. 이를 학습용 절반에서 구한 통합 평균과 범주별 평균 모두와 비교하며, 각각의 제곱 크기로 정규화한다. 난수 시드는 2026091201이다. 표집 과정에서 같은 홀드아웃 작품이 반복될 수는 있지만, 해당 분할의 참조 대상 추정에는 어떤 홀드아웃 작품도 사용하지 않는다. 개발용 스케일러는 계속 고정한다.

소수 표집의 영향과 비교 집합 변화의 영향을 구별하기 위해, 학습 참조 대상과 정규화 인자를 고정한 채 각 분할의 정확한 홀드아웃 평균으로 기대 점수도 계산한다(표~\ref{tab:painting-control-ranges}).

\begin{table}[H]
\centering\small
\begin{tabular}{@{}llcc@{}}
\toprule
대조 표집 & 참조 대상 & 두 번 추출 & 정확한 홀드아웃 평균\\\midrule
통합 & 통합 & $[-.758,1.380]$ & $[.135,.428]$\\
범주 층화 & 통합 & $[-.554,2.400]$ & $[.544,.993]$\\
범주 층화 & 범주별 & $[-.248,1.946]$ & $[.517,.954]$\\
\bottomrule
\end{tabular}
\caption{분할 1,000회에 걸친 실제 회화 대조 점수의 중앙 95\% 범위. 정확한 평균은 각 분할 내 표집을 적분하여 제거한다. 이 유한 집합 범위는 신뢰구간이 아니다.}
\label{tab:painting-control-ranges}
\end{table}

통합 대조 분석에서는 표집 변동을 적분해 제거한 뒤 표준편차가 $.561$에서 $.076$으로 낮아진다. 이러한 대조 분석도 여전히 불완전한 제목 범주와 작품당 단일 출처 촬영본에 조건부이며, 예술적 충실도를 보정하는 것은 아니다.

\begin{table}[H]
\centering\small
\begin{tabular}{@{}lrrrr@{}}
\toprule
 & \multicolumn{2}{c}{공통 장면 11개} & \multicolumn{2}{c}{장면 14개}\\
모델 & 통합 대상 & 범주별 대상 & 특징군 균등 & 공분산\\\midrule
GPT Image 1 & 1.557 & 1.366 & 1.551 & 2.103\\
GPT Image 2 & 1.099 & 1.130 & 1.243 & 1.511\\
2.5 Flare & 1.726 & 1.539 & 1.765 & 1.636\\
2.5 Sunburst & 1.633 & 1.525 & 1.680 & 1.617\\
Nano Banana 2 & 0.987 & 0.948 & 1.203 & 1.385\\
FLUX.2 Max & 0.750 & 0.786 & 0.808 & 0.928\\
\bottomrule
\end{tabular}
\caption{대안적인 참조 대상과 지표에서의 기술적 오차. 범주별 대상은 제목에서 도출한 수면·건축·육지 층을 사용하며 혼합 장면 세 개는 제외한다. 각 열은 고유한 참조 정규화 인자를 사용하므로 열 간 절댓값을 직접 비교할 수 없다. 공분산은 개발 작품만으로 적합하며, 고정된 50\% 수축을 적용한다.}
\label{tab:review-targets}
\end{table}
\FloatBarrier

\subsection{크기, 화가 포괄 범위, 좌표 가중치}\label{app:magnitude-coverage}
식~\ref{eq:scene-decomposition}의 반복 간 버전은 보관된 벡터에 대해 정확히 성립한다. 그 변동항은 제곱 이탈을 장면 중심화된 반복들 사이의 교차곱으로 대체한다. 독립 반복에 대한 모집단 대상은 음이 아니지만, 유한 표본에서 이 항은 음수일 수 있다. 스칼라 배율 보정은 다른 장면 13개만으로 적합한다. 모든 폴드에서 분모 $Q$는 양수이다. GPT Image 2에 적합한 스칼라는 .438--.459, FLUX는 .598--.694이며, 폴드별 모든 점수를 보관한다. 어떤 보정 매개변수도 자신의 홀드아웃 점수를 최소화하는 방식으로 선택하지 않는다. 이후 삭제 점검에서는 겹치는 폴드를 독립적인 재현 실험으로 취급하지 않고, 남은 장면 13개의 각 패널에서 전체 절차를 다시 적합한다. 그 평균 GPT Image 2--FLUX 오차 차이는 $-.189$에서 $-.137$ 사이이다.

\begin{table}[H]
\centering\small
\begin{tabular}{@{}lrr@{}}
\toprule
모델 & 대조값의 제곱 크기 $Q$ & 정렬 비 $\beta/\sqrt Q$\\\midrule
GPT Image 1 & 2.449 & 0.600\\
GPT Image 2 & 2.223 & 0.670\\
2.5 Flare & 2.281 & 0.511\\
2.5 Sunburst & 2.289 & 0.545\\
Nano Banana 2 & 0.994 & 0.444\\
FLUX.2 Max & 0.741 & 0.546\\
\bottomrule
\end{tabular}
\caption{화가 대조값의 크기와 정렬에 대한 기술적 요약. $Q=1$은 참조의 제곱 크기와 같으며, $D=1-2\beta+Q$이다. 유한 표본에서 정렬 비는 $[-1,1]$ 밖에 있을 수 있다.}
\label{tab:review-alignment}
\end{table}

화가 쌍을 비교할 때 생성·참조 대조값은 모두 두 화가 간의 직접 차이이며, 분모는 참조 쌍 거리의 제곱이다. 한 화가를 제외할 때는 나머지 세 화가의 생성·참조 조건을 다시 중심화한다. 참조 특이 성분은 중심화된 $4\times31$ 참조 행렬의 랭크 1 항 세 개이다. 각 성분은 특정 화가 하나를 나타내는 것이 아니라 화가 대조와 특징 방향을 결합한다. 성분별 정렬 반응은 각 항에 별도로 투영한다.

\begin{table}[H]
\centering\small\setlength{\tabcolsep}{4pt}
\begin{tabular}{@{}lrrr|rrrr@{}}
\toprule
 & \multicolumn{3}{c}{성분별 정렬 반응} & \multicolumn{4}{c}{해당 화가 제외 후 정렬 반응}\\
모델 & 제1 & 제2 & 제3 & 모네 & 시슬레 & 피사로 & 세잔\\\midrule
GPT Image 1 & 1.248 & 0.310 & 0.356 & 1.140 & 1.072 & 0.840 & 0.389\\
GPT Image 2 & 1.156 & 0.617 & 0.800 & 1.195 & 0.972 & 0.992 & 0.651\\
2.5 Flare & 1.022 & 0.211 & 0.383 & 1.034 & 0.734 & 0.802 & 0.222\\
2.5 Sunburst & 1.082 & 0.250 & 0.417 & 1.116 & 0.840 & 0.762 & 0.284\\
Nano Banana 2 & 0.513 & 0.371 & 0.195 & 0.569 & 0.446 & 0.388 & 0.276\\
FLUX.2 Max & 0.666 & 0.021 & 0.183 & 0.539 & 0.511 & 0.527 & 0.096\\
\bottomrule
\end{tabular}
\caption{정렬 반응이 화가들을 얼마나 포괄하는지에 대한 기술적 요약. 참조 성분은 중심점 변동의 66.3\%, 20.6\%, 13.0\%를 설명한다. 화가를 하나 제외할 때마다 참조 대상을 다시 계산하고 나머지 세 화가를 다시 중심화한다.}
\label{tab:review-components}
\end{table}

표~\ref{tab:review-components}는 전체적으로 양의 정렬이 나타나더라도 그 정렬이 각 성분에 고르게 분포하지 않음을 보여준다.

특징군 균등 가중치는 각 좌표에 해당 특징군 크기(11, 8, 12)의 역수를 부여한다. 공분산 민감도 분석에는 화가별 총가중치를 같게 한 개발 작품 221점을 사용한다. 그 가중 공분산을 $\Sigma$라고 하면, 지표는 $[.5\Sigma+.5\operatorname{tr}(\Sigma)I/31]^{-1}$이다. 수축 정도는 고정하며 평가 이미지에 맞추어 조정하지 않는다. 이러한 가중치 선택이 독립적으로 검증된 예술적 지표인 것은 아니다.

좌표를 하나씩 제외하는 31개 경우 모두에서 FLUX는 가장 낮은 점 추정치를 유지한다. 보고서는 음의 항을 포함해 모든 좌표의 기여도를 보관한다. 텍스처를 제외해도 Flare와 Sunburst의 오차 추정치는 각각 1.580과 1.604로, 대조값이 없는 기준값 1보다 높다.

\subsection{설계를 반영한 에너지 교정}
주 에너지 값에는 대각선의 0 거리까지 포함한 경험분포를 사용한다. 고정된 경험적 참조 분포와 장면 조건부 출력 분포의 동일 가중 혼합 사이 에너지를 추정할 때, 서로 다른 장면의 생성 쌍에는 이미 필요한 독립 추출이 갖추어져 있다. 같은 장면의 생성 항만 조정하면 된다. 반복 두 번과 장면 $S$개일 때 기술적 교정은 다음과 같다.
\[
 \widehat{\mathcal E}_{\rm strata}
 =\widehat{\mathcal E}_{\rm empirical}
 -\frac1{2S^2}\sum_s\|y_{s1}-y_{s2}\|.
\]
이 평가 대상은 유한한 참조 패널에 조건부이므로 참조 내 거리는 경험값으로 유지한다. 통합 IID 배율 $n/(n-1)$을 적용하면 서로 다른 장면 간 항까지 달라져 다른 대상을 추정하게 된다. 이 교정 후에도 화가 지정에서 일반 유화 조건을 뺀 차이 21개가 음수라는 결과는 유지된다. 하지만 역사적·생성 이미지의 내용 구성이 같지 않다는 문제는 해결하지 못하며, 새로운 에너지 검정도 수행하지 않는다.

\subsection{불확실성과 합성 데이터 포함률 점검}
Student 절차는 장면 요약값 14개를 사용한다. 고정된 조건부 평균 $\delta_s$를 갖는 독립 장면 오차의 경우, 추정 분산의 기댓값은 다음과 같다.
\[
 \mathbb E[s_\delta^2/S]
 =\operatorname{Var}(\bar{\widehat\delta})
 +\frac1{S(S-1)}\sum_s(\delta_s-\bar\delta)^2.
\]
따라서 요청 잡음뿐 아니라 실제 장면 간 이질성도 포함된다. 이것이 해당 고정 패널에서 새로 요청할 때의 분산을 직접 추정하는 값이 아닌 이유이다. 정확한 포함률을 보장하지도 않는다. 비정규적인 교차곱과 의존성이 Student 근사에 영향을 줄 수 있기 때문이다.

초기의 5,000회 실행 시나리오(시드 2026091202)는 여섯 구성, 장면 14개, 반복 두 번, 평가 항목 21개의 집합을 사용하고, 참조 크기를 1로 둔다. 가우시안 오차, 분산을 표준화한 $t_3$/이분산 오차, 고정 장면 상호작용 오차의 동시 포함률은 각각 95.96\%, 97.42\%, 99.68\%이다. 이들의 독립적인 중심화 잡음 파워는 $.5$이다. 잡음의 절반을 모든 장면과 두 반복이 공유하는 모델별 상태에 할당하면 포함률이 0.04\%로 낮아지고, 교정 오차의 평균 편향은 약 $.25$가 된다.

후속 탐색에서는 파워 $.5$, $1.5$, $3$의 독립 가우시안 오차를 사용하고, 공분산을 등방성 또는 참조 성분과 정렬된 랭크 3으로 설정한다(경우당 5,000회 실행). 포함률은 95.06--96.80\%이며, 몬테카를로 표준오차는 $.25$--$.31$퍼센트포인트이다. 비교를 위해 반복 차이를 직접 중심화해 추정한 관찰 파워는 $H$ 대비 $.424$--$2.596$이다. 이 스트레스 모형은 포함률 근사를 점검하는 것이지 실제 서비스 독립성을 검증하는 것이 아니다. 이를 바탕으로 유의성 임계값을 선택하지도 않는다.

\subsection{출처 품질 민감도 분석}\label{app:reference-quality}
AI 보조 도구가 콘택트 시트와 주변부 확대 영역을 이용해 참조 복제 이미지 649개와 개발용 복제 이미지 221개를 모두 검사하고, 제거 가능한 보정 표적, 프레임, 여백, 라벨을 기록한다. 경계가 불확실하면 변경하지 않는다. 참조 내용은 제목 라벨과 비교하기 전에 시각적으로 분류하며, 혼합되거나 불명확한 경우에는 제목 라벨을 유지한다.

감사 결과 참조 이미지 90개와 개발 이미지 41개에서 크롭을 선정했으며, 그중 색 보정 띠 사례는 열 개이다. 불확실한 경계 43개는 변경하지 않았다. 시각 범주는 제목 기반 할당 230개와 달랐고, 혼합되거나 불명확한 사례 42개는 원래 라벨을 유지했다. 이러한 AI 코딩 민감도 분석은 사람의 검증을 거치지 않았으며, 독립적인 정답 자료가 아니다.

감사된 직사각형 크롭은 원래의 방향·ICC 규칙을 따른다. 그다음 종횡비를 유지하면서 짧은 변을 512픽셀로 맞추는 Lanczos 리사이징을 업샘플링 없이 적용한다. 특징 31개와 모든 작품을 유지한다. 먼저 교정한 참조 자료를 원래 스케일링으로 비교한 뒤, 개발용 IQR 스케일링을 다시 적합하고 생성 벡터도 일관되게 변환한다. 라벨 민감도 분석에는 동일한 장면 11개를 사용한다. 버전 기록에는 감사 결정, 새 벡터, 원래 결과를 보존하며, 모델 점수에 따라 교정을 선택하지 않는다.

어느 스케일러를 사용하든 여섯 정렬 반응 구간은 모두 0보다 높게 유지된다. 원래 스케일링으로 교정 영역을 사용하면 FLUX의 $D=0.832$이다. 다시 적합한 뒤의 (보정 전, 홀드아웃 배율 보정 후) 오차는 GPT Image 1 $(1.631,0.693)$, GPT Image 2 $(1.259,0.579)$, Flare $(1.737,0.760)$, Sunburst $(1.577,0.711)$, Nano Banana 2 $(1.121,0.846)$, FLUX $(0.872,0.762)$이다. 각각의 최솟값은 여전히 FLUX와 GPT Image 2가 갖는다.

\section{이미지 선정과 출처}\label{app:images}
생성 예시는 사전 선언한 첫 번째 장면과 첫 번째 반복을 사용한다. 분석 후 선정했지만, 외관이나 점수에 따른 선별은 하지 않았다. 장면 설명은 다음과 같다.
\begin{quote}
‘넓은 강이 자갈 강둑을 따라 굽이친다. 맞은편 강둑에는 키 큰 나무 세 그루가 서 있고, 그 뒤로 낮은 언덕이 있으며, 하늘은 옅게 흐려 있다.’

\textit{A broad river bends past a gravel bank. Three tall trees stand on the far bank, with a low hill behind them under a pale overcast sky.}
\end{quote}
모든 프롬프트는 ‘문자나 프레임을 넣지 말 것’(No text or frame.)으로 끝난다. \ref{sec:design}절의 지시문만 달라진다.

그림~\ref{fig:original-generated}의 참조 작품은 완료된 감사에서 수면 중심으로 분류되고 퍼블릭 도메인 상태가 기록된 측정 작품 중, 각 화가별 작품 ID의 사전식 순서가 가장 앞선 작품이다. 시슬레·피사로 복제 이미지에는 색 보정 띠를 제거하는 감사된 크롭을 사용했고, 모네·세잔은 전체 영역을 유지했다. 매니페스트에는 이미지 식별 정보, 출처, 크롭 상자, 정확한 프롬프트가 기록되어 있다. 이 작품들은 생성 장면 설명에 맞춘 것이 아니다.
\FloatBarrier

\section{전체 모델 및 화가 요약}\label{app:results}
표~\ref{tab:all-specificity-pairs}--\ref{tab:artist-distributions}와 그림~\ref{fig:specificity-diagnostics}--\ref{fig:specificity-distributions}는 모델 비교, 오차 성분, 집합 수준의 근접도와 산포를 요약한다. \emph{에너지 불일치}는 집합 간 거리와 내부 산포를 결합하며, 값이 낮을수록 더 가깝다(부록~\ref{app:estimation}, \ref{app:diagnostics}). 화가 지정은 일반 유화 지시에 비해 모델--화가 조합 24개 중 21개에서 에너지를 낮추며, 반복 장면 교정을 적용해도 같다. 화가를 지정한 집합 24개 모두에서 전체 특징 분산이 대응하는 역사적 작품보다 작다. 이 기술적 비교는 이름 지정 또는 일반 유화 출력 28개를 기록된 참조 패널과 비교하며, 소재 구성이 서로 다르다. 따라서 화가 대조값의 일치를 입증하는 것은 아니다.

\begin{table}[H]
\centering\small
\begin{tabular}{@{}llr@{}}
\toprule
모델 A & 모델 B & $D_A-D_B$ [동시 구간]\\\midrule
GPT Image 1 & GPT Image 2 & $0.346\;[-0.400,1.093]$\\
GPT Image 1 & 2.5 Flare & $-0.165\;[-1.027,0.697]$\\
GPT Image 1 & 2.5 Sunburst & $-0.069\;[-1.181,1.043]$\\
GPT Image 1 & Nano Banana 2 & $0.463\;[-0.602,1.528]$\\
GPT Image 1 & FLUX.2 Max & $0.771\;[-0.037,1.580]$\\
GPT Image 2 & 2.5 Flare & $-0.512\;[-1.150,0.126]$\\
GPT Image 2 & 2.5 Sunburst & $-0.415\;[-1.181,0.350]$\\
GPT Image 2 & Nano Banana 2 & $0.116\;[-0.922,1.155]$\\
GPT Image 2 & FLUX.2 Max & $0.425\;[-0.161,1.011]$\\
2.5 Flare & 2.5 Sunburst & $0.096\;[-0.286,0.478]$\\
2.5 Flare & Nano Banana 2 & $0.628\;[-0.283,1.540]$\\
2.5 Flare & FLUX.2 Max & $0.937\;[0.256,1.618]$\\
2.5 Sunburst & Nano Banana 2 & $0.532\;[-0.436,1.500]$\\
2.5 Sunburst & FLUX.2 Max & $0.840\;[0.058,1.623]$\\
Nano Banana 2 & FLUX.2 Max & $0.309\;[-0.847,1.464]$\\
\bottomrule
\end{tabular}
\caption{사전 명시한 추론 비교 집합에 속하는 대응 모델 비교 15개 전체. 구간은 평가 항목 21개 전체에 대해 보정한 근사 95\% 동시 구간이다. 음의 차이는 명시한 특징 기하 구조 평가 대상에서 모델 A에 유리함을 뜻한다.}
\label{tab:all-specificity-pairs}
\end{table}

\begin{table}[H]
\centering\small
\begin{tabular}{@{}llrrrr@{}}
\toprule
모델 & 측정값 & 모네 & 시슬레 & 피사로 & 세잔\\\midrule
GPT Image 1 & 에너지 변화 & $-2.515$ & $-2.739$ & $-2.628$ & $-1.336$\\
 & 분산비 & $0.479$ & $0.694$ & $0.603$ & $0.598$\\
\addlinespace[3pt]
GPT Image 2 & 에너지 변화 & $-0.602$ & $-1.092$ & $-0.195$ & $-0.373$\\
 & 분산비 & $0.283$ & $0.312$ & $0.303$ & $0.303$\\
\addlinespace[3pt]
2.5 Flare & 에너지 변화 & $0.962$ & $-0.008$ & $0.429$ & $-1.013$\\
 & 분산비 & $0.269$ & $0.371$ & $0.321$ & $0.275$\\
\addlinespace[3pt]
2.5 Sunburst & 에너지 변화 & $0.267$ & $-0.718$ & $-0.091$ & $-1.653$\\
 & 분산비 & $0.248$ & $0.274$ & $0.273$ & $0.318$\\
\addlinespace[3pt]
Nano Banana 2 & 에너지 변화 & $-0.513$ & $-1.381$ & $-0.741$ & $-1.102$\\
 & 분산비 & $0.829$ & $0.708$ & $0.572$ & $0.877$\\
\addlinespace[3pt]
FLUX.2 Max & 에너지 변화 & $-0.707$ & $-0.886$ & $-1.080$ & $-1.431$\\
 & 분산비 & $0.340$ & $0.416$ & $0.570$ & $0.622$\\
\bottomrule
\end{tabular}
\caption{보완적인 경험분포 요약. 에너지 변화는 화가 지정 조건의 참조 에너지에서 일반 유화 조건의 참조 에너지를 뺀 값이며, 분산비는 생성 분산을 참조 분산으로 나눈 값이다. 음의 변화는 화가 지정에 유리하고, 1보다 작은 비율은 생성 특징 분포가 더 좁음을 나타낸다.}
\label{tab:artist-distributions}
\end{table}
\FloatBarrier

\begin{figure}[H]
\centering\sourcefigure{specificity_diagnostics.pdf}
\caption{원래 참조 대상에 대한 기술적 요약. (a) 진폭 오차와 직교 잔차, 즉 전체 라벨별 참조 패턴에 수직인 차이를 합하면 $D$가 된다(식~\ref{eq:decomposition}). (b) 모든 모델·화가의 생성/참조 분산비. 전체 산포에는 고정된 장면 설명 간 변동이 포함되며, 모드 붕괴 지표가 아니다.}
\label{fig:specificity-diagnostics}
\end{figure}

\begin{figure}[H]
\centering\sourcefigure{specificity_sunburst_distributions.pdf}
\caption{GPT Image 2.5 Sunburst의 화가 지정 출력은 네 투영 모두에서 참조 작품보다 더 집중되어 보인다. 이 모델 예시는 주 특징 결과를 확인하기 전에 선정했다. 각 패널은 해당 화가의 참조 작품만으로 적합한 주성분을 사용하므로 화가마다 축이 다르다. 생성 장면 설명 14개와 두 반복의 내용 분포는 역사적 패널과 다르다. 투영은 기술적 시각화이며, 수치 요약에는 좌표 31개 전체를 사용한다.}
\label{fig:specificity-distributions}
\end{figure}
\FloatBarrier

\section{보완적 개입 실험}\label{app:cohorts}
이전의 자료 집합들은 관련되지만 서로 다른 질문을 다룬다(표~\ref{tab:cohorts}). 여섯 모델 실험과 통합하지 않으며, 독립 연구자가 수행한 재현 실험으로 취급하지도 않는다. 과거의 모델 라벨은 검증된 내부 모델이 아니라 기록된 요청 구성을 식별한다(부록~\ref{app:provenance}).

\begin{table}[H]
\centering\small
\begin{tabularx}{\linewidth}{@{}>{\raggedright\arraybackslash}p{.21\linewidth}>{\raggedright\arraybackslash}p{.30\linewidth}>{\raggedright\arraybackslash}X@{}}
\toprule
자료 집합 & 비교 & 본 논문에서의 역할\\\midrule
네 화가 탐색 & 참조 649개, 생성 출력 1,920개 & 참조·생성 분리를 기술하고 통제 실험의 동기를 제공한다.\\[5pt]
통제된 이름 지정 & 출력 1,006개, 모네 참조 38개·세잔 참조 32개 & 이름 지정/미지정 비교 여섯 개, 상세/짧은 프롬프트 비교 두 개, 후속 변환 및 조건부 분산 진단.\\[5pt]
팔레트 개입 & 서로 별개인 출력 192개 집합 두 개 & 이름 지정--일반 유화 간 채도 반응; 집합을 통합하지 않는다.\\[5pt]
일반 지시문 비교 & 새 장면 24개의 세잔/일반 유화 출력 96개 & 일반 회화 지시를 넘어서는 근접도 향상을 검증한다.\\[5pt]
고정 매핑 전이 & 새 장면 12개의 FLUX/세잔 출력 240개 & 과거의 에너지·예측 오차 순서가 서로 반대였던 현상에 대한 전향적 반례.\\
\bottomrule
\end{tabularx}
\caption{보완 근거는 새로운 여섯 모델 실험과 분리해 유지한다. 서로 다른 프롬프트·패널·유한 표본 평가 대상을 하나의 통합 검정으로 합치지 않는다.}
\label{tab:cohorts}
\end{table}

\subsection{이름 지정에 따른 근접도와 산포}
네 화가 탐색에서는 화가 지정 출력 1,536개, 화가 미지정 출력 384개, 프롬프트 방식 세 개, 장면 16개를 사용했다. 모델--화가--방식 조합 24개의 전체 분산비는 .206--.376이고, 그룹을 고려한 생성/참조 판별 정확도는 .940--.983이었다. 모네·시슬레·피사로·세잔의 통합 화가 지정 에너지 중앙값은 각각 2.068, 1.678, 1.975, 2.142였으며, 대응하는 화가 미지정 중앙값은 1.732, 1.878, 1.854, 1.657이었다. 통합 화가 지정 중앙값이 더 낮은 화가는 시슬레뿐이었다. 크기를 맞춘 실제/실제 비교의 중앙값은 .193, .215, .164, .199였다. 이러한 사후 요약은 해당 탐색을 기술할 뿐, 일반적인 이름 지정 효과를 뜻하지 않는다.

후속 통제 패널에서는 탐색 이후 모네와 세잔을 선택했다. 장면 24개에 모델 세 개, 지시문 두 개, 반복 세 번을 교차 적용했으며, 짧은 장면 출력 142개로 프롬프트 길이 비교 두 개를 추가했다. 참조 패널은 모네 38점, 세잔 32점으로 이루어졌고, 수면·건축·육지 비중은 관리자가 실행한 LLM 시각 코딩에 근거했다. 이름 지정에서 미지정을 뺀 에너지 변화는 Nano Banana 2에서 $-.846,-1.818$, FLUX에서 $-.625,-.896$, GPT Image 2에서 $-.234,-.056$이었다(모네, 세잔 순). 처음 네 귀무가설은 원래 여덟 개 검정 집합에서 기각되었다. 별도의 FLUX 출력 72개 집합에서도 두 방향을 재현했으며($-.695,-.952$), 그 자체의 네 평가 항목 집합에서 기각되었다.

새 장면 24개에 대한 독립적인 GPT Image 2/세잔 출력 96개 비교에서는 일반 지시문 대비 에너지가 1.195 감소했고($p=.00001$, 고정 $\alpha=.025$), 화가 지정/일반 유화 전체 분산비는 .552였다. 두 위치로 구성된 블록 48개 안에서 지시문을 무작위 배정하고, 참조 범주 비중 $(3,11,18)/32$를 사용했다. 블록 내 부호 추출 99,999회와 플러스원 교정을 이용해 모든 개별 효과가 0이라는 강한 귀무가설(sharp null)을 검정했다. 이는 지시문 효과를 검정하는 것이지 에너지의 동일성만 검정하는 것은 아니다. 따라서 해당 패널에서 일반 지시문은 근접도 향상 전체를 재현하지 못했지만, 근접도가 올바른 네 화가 대조값을 입증하지는 않는다. 마찬가지로 전체 산포가 작다고 해서 장면 정보가 손실되었다고 볼 수는 없다. 통제된 FLUX/모네 비교에서 반복을 홀드아웃한 장면 검색 정확도는 이름 지정 시 44.4\%에서 59.7\%로 높아졌다. 이것은 상대적인 특징 구별 가능성을 측정하는 것이지, 검증된 프롬프트 준수도를 측정하는 것은 아니다.

\subsection{팔레트 반응성과 반대 방향의 전이 결과}
두 팔레트 자료 집합은 각각 출력 192개로 구성되며, 과거 GPT Image 2 별칭을 요청했다. 장면 여섯 개, 반복 네 번, 차분한/선명한 색 지시, 화가 미지정·일반 유화·모네·세잔 지시문을 사용했다. 주 분석의 이름 지정--일반 유화 채도 반응 상호작용은 통계적으로 구별되지 않았다. 모네·세잔 순으로 첫 집합에서는 $-.284,-.018$, 후속 집합에서는 $-.159,.037$이었다. 보정된 구간은 원래의 두 상호작용 집합과 후속 네 평가 항목 집합 모두에서 0을 포함했으며, 후자에는 FLUX 이름 지정 대조 두 개도 들어 있었다. 부차적인 일반 유화--화가 미지정 반응은 $-.853$과 $-.891$이었고, 명목 95\% 구간이 0을 제외했다. 따라서 추가적인 이름 효과가 확인되지 않더라도 일반 회화 지시가 측정된 팔레트 반응을 약화할 수 있다. 팔레트 수준이 두 개뿐이므로 이득 감소, 포화, 작동 범위의 이동을 구분할 수 없으며, 상호작용이 구별되지 않았다는 사실이 동등성을 입증하지는 않는다.

이전 매핑은 화가 미지정 특징 벡터에 평행이동을 적용하고, 선택적으로 등방성 스케일링을 더하여 화가 지정 프롬프트의 장면 평균을 예측했다. 생성 이미지만으로 적합하고 홀드아웃 장면에서 평가했다. 모델--화가 평균 여섯 개 모두에서 실제 이름 지정의 참조 에너지가 평행이동과 등방성 스케일링 조합보다 낮았다. 스케일링을 추가하면 여섯 개 모두에서 화가 지정 평균의 예측은 개선되지만, 다섯 개에서는 참조 에너지가 나빠졌다. 그러나 원래의 적합을 그대로 사용해 새 장면 12개와 반복 열 번으로 수행한 전향적 FLUX/세잔 출력 240개 비교에서는 이러한 순서가 전이되지 않았다. 평행이동만 적용한 경우에 비해 평행이동/스케일링은 에너지와 반복 교정 예측 오차를 모두 개선했다. 차이는 각각 $-.390\;[-.509,-.271]$과 $-5.047\;[-7.880,-2.214]$이며, 근사 동시 잭나이프 구간을 사용했다. 실제로 얻은 예측 오차 구간의 반폭은 계획 기준을 초과했다. 이 반대 결과는 추가 표집 없이 그대로 유지한다. 따라서 과거의 상반된 순서를 보편적인 메커니즘으로 취급할 수는 없다. 이러한 화가 지정/미지정 매핑은 표~\ref{tab:review-calibration}의 화가 대조값에 대한 사후 배율 보정과 다르다.

\clearpage
\section{수집 이력과 계산 범위}\label{app:provenance}
과거의 GPT Image 1·GPT Image 2 라벨은 요청한 모델 이름을 식별한다. 이후 수행한 기술적 음성 대조 실험에서는 이전 이미지 인터페이스가 임의로 만든 이미지 모델 식별자도 받아들인다는 사실이 확인되었다. 이것이 유효한 이름들이 동일한 내부 모델을 선택한다는 뜻은 아니지만, 요청이 성공적으로 수락되었다는 사실만으로 두 모델의 구분을 인증할 수는 없다. 따라서 과거의 대조는 기록된 요청 그룹 간 비교로 해석한다. 현재 실험은 명시적이고 문서화된 모델 식별자와 고정 경로를 사용하며, 과거 이미지를 새로운 변형 모델의 검증된 예시로 재사용하지 않는다. 현재 출력은 모두 요청한 정사각형 크기로 디코딩되지만, 응답 본문에는 모델·품질 필드가 다시 표시되지 않는다. 따라서 라벨은 독립적인 체크포인트 인증 없이, 문서화된 요청 구성을 나타낸다. 내부 체크포인트와 학습 코퍼스는 검사할 수 없다.

이번 실험에서는 기술적 점검과, 모델 선택의 적격성을 확인하지 못해 특징 측정 전에 종료한 이전 시도를 제외했다. 이 자료 집합들 사이에 공유하는 이미지는 없다. 마찬가지로 이전의 불완전한 일반 지시문 집합도 완전한 후속 이미지 96개 집합과 통합하지 않는다. 정확한 배정, 시도 결과, 측정 제외 내역은 버전별 매니페스트에 보관되어 있으며, 시각적 성공 여부를 기준으로 선정하지 않았다.

\end{document}
